% 本文件由 LaTeX 排版 Agent 根据有效 translation.json 自动生成。
% author=Eusebius of Caesarea (c. 265-c. 340); work_id=2504; title=赞颂君士坦丁演说

\chapter{赞颂君士坦丁演说}

% structure: p0001 source=h1 target=omitted reason=page-title-already-rendered-by-parent

\Emph{在他登基三十周年纪念日发表。}\par

% structure: p0003 source=h2 target=section reason=relative-heading-rank; base=h2

\section{序言}

一、我不是来呈现虚构的叙述，也不是以优雅的语言来取悦耳朵，仿佛想用塞壬的声音迷惑听众；我也不是把装在饰有可爱花朵的金杯中的快乐之饮（我是指风格的优雅）献给那些喜爱此类事物的人。我宁愿遵循智者的训诫，劝告所有人避开并远离那条常走的路，使自己不接触粗俗的群众。\par

二、那么，我准备好了用更新的格调来颂扬我们的皇帝；虽然渴望与我同担当前任务的人数不胜数，我决心避开常人的道路，追求那条未经践踏的小径，没有洗净双脚的人不得踏足。让那些崇拜粗俗风格、充满幼稚精巧、并讨好愉快而流行的缪斯的人，既然快乐是他们追求的目标，就试图通过叙述仅仅人类的功绩来取悦人的耳朵。但是，那些通晓普遍科学、既达到神圣知识又达到人类知识、并认为选择后者才是真正卓越的人，将更喜欢皇帝那些上天本身赞许的德性，以及他虔诚的行为，而不是他仅仅人类的成就；并将把歌颂他较小功绩的任务留给较低级的颂扬者。\par

三、既然我们的皇帝既拥有那与天主直接相关的神圣智慧，又拥有涉及人类利益的知识；那么，让那些胜任此任务的人描述他世俗的成就——这些成就尽管伟大而超越，并有益于人类（因为皇帝的一切都是伟大而高尚的），但相对于他神圣的品格而言，对站在神圣庭院之外的人来说仍然是次要的。\par

四、但是，让那些在圣所之内、能进入其最内层和未经践踏的隐秘处的人，对每一只不洁的耳朵关上大门，仿佛只对初入门者展示皇帝品格的秘密奥迹。让那些在虔敬的溪流中净化了耳朵、并用心灵本身的翱翔之翼提升思想的人，加入环绕万有之主的群体，并沉默地学习神圣的奥迹。\par

五、同时，让那些神圣的神谕成为我们在这些奥迹中的导师——这些神谕不是由占卜之灵（或者更确切地说，是疯狂和愚蠢之灵）所赐，而是由神圣真理的默感所赐；它们一般地谈论主权：论及万有的至高主权者，以及环绕万有之主的天庭阵列；论及我们面前那个帝国权能的模范，以及那个伪币；最后，论及两者产生的后果。因此，既然这些神谕在我们对神圣礼典的认识中启蒙我们，让我们如下开始我们的神圣奥迹。\par

% structure: p0009 source=h2 target=section reason=relative-heading-rank; base=h2

\section{第一章}

一、今天是我们的伟大皇帝的庆节：我们他的子女因此欢喜，感受到神圣主题的默感。主持我们庆典的是那位伟大的君主本身；我是指那真正伟大的那一位；关于他，我断言（那位听到我的皇帝也不会生气，反而会赞成这种对天主的赞美）：他高于并超越一切受造物，是最崇高者、最伟大者、最全能者；他的宝座是穹苍，大地是他脚下的脚凳。他的存有没有人能恰当地理解；环绕他的荣耀的不可言喻的光辉，从每一只眼睛前驱散凝视他神圣威严的目光。\par

二、他的臣仆是天军；他的军队是上界的能力，他们承认他是他们的主人、主和君王。不可计数的天使群体、总领天使的团队、圣洁神灵的合唱团，从永远光明的源泉汲取并反射他的光辉。是的，每一光明，特别是那些位于天界之外的神圣无形理智，都以崇高神圣的赞歌庆祝这位庄严的君主。广阔的苍天，像一层面纱，介于外面的人和居住在他王宫中的人之间；在这苍天周围，太阳、月亮以及其余天上的光体（如同皇宫入口周围的执炬者），为了他们的君主，执行着指定的轨道；按照他命令的话语，向那些命运落在天界之外较黑暗区域的人，持着永不熄灭的光。\par

三、确实，当我记起我们自己的得胜皇帝向这位全能君主献上赞颂时，我很好地跟随他，因为我知道，唯独从他，我们才拥有我们赖以生存的帝国权能。虔诚的凯撒们，受他们父亲智慧的教导，承认他是每一样祝福的源泉：军人、全体人民，无论是在帝国乡村还是城市中，连同各省总督，按照他们伟大的救主和导师的训导聚集在一起，崇拜他。简言之，全人类家庭，每个民族、支派和语言，集体和分别地，尽管在其他问题上观点多样，在这一个信仰宣认上却是一致的；并且，由于他们内植的理性和他们心灵自发的、未经教导的冲动，他们联合呼求独一的天主。\par

4. 难道大地这普世的框架不承认他是她的主，并且藉着她所出产的植物和动物生命宣示她服从更高权能的旨意吗？川流不息的江河，从隐秘无尽的深处涌出的永恒泉水，将她们奇妙之源的成因归于他。大海的浩瀚之水，被封闭在深不可测的深渊室中，那高涌的波涛，似乎威胁着大地本身，但一接近岸边就退缩，被祂神圣律法的力量所遏制。冬雨的适时降落，滚滚雷声，闪电的闪光，风的漩涡气流，云的空中轨迹，这一切都向那些祂的位格看不见的人启示祂的同在。\par

5. 那极光明的太阳，历经岁月永不停歇地运行，只承认他为主，顺从祂的旨意，不敢偏离指定的路径。月亮的次等光辉，按期交替盈亏，服从于祂的神圣命令。苍天的美丽机制，闪耀着众星，以和谐秩序运行，保持每条轨道的尺度，宣告祂是一切光的赐予者。是的，所有天上的发光体，都按照祂的旨意和话语，保持伟大而完美的运动统一，追寻埃塞俄比亚轨道，在流转的岁月中完成它们遥远的行程。昼夜的交替，季节的更替，宇宙的秩序与比例，都宣告祂[祂无限的力量]的多重智慧。对祂而言，那在空间伸展中运行的无形德能，献上应有的赞美颂歌。对这个地球本身，对上面的诸天，对天穹之外的歌队，都向祂伟大的君主致敬：天使的军旅以不可言喻的赞美诗歌迎接祂；从非受造之光中取得存在的诸灵，敬拜祂为他们的造物主。在天地之前的永恒世代，连同其他无限的时期，先于一切可见受造，都承认祂是唯一至高无上的主和统治者。\par

6. 最后，那位在万有之中、在万有之前和万有之后者，祂的独生、先在的\OriginalTerm{圣言}{Word}，伟大天主的\OriginalTerm{大司祭}{Great High Priest}，比一切时间和世代更古老，专一为祂圣父的荣耀，首先且独自为人类的获救向祂代祷。宇宙至高无上统治者，同享祂圣父王国的光荣：因为祂就是\OriginalTerm{光}{Light}，那超越宇宙，环绕圣父的位格，介入并划分永恒非受造本质与一切受造存在之间的光；那从上流下，源自不知起源和终结的天主，照亮超天界区域和天本身所包含的一切，以智慧的光辉明亮远超太阳的光辉。祂是这整个世界拥有至高统治权的，祂在万有之上和万有之内，贯穿一切可见与不可见者；天主的圣言。我们深受天主宠爱的皇帝，从祂并由祂，好似接受了一份神圣统治的抄本，效法天主自己，指导这个世界事务的行政管理。\par

% structure: p0016 source=h2 target=section reason=relative-heading-rank; base=h2

\section{第二章}

1. 天主的独生\OriginalTerm{圣言}{Word}，从无始的世代直到无限无尽的世代，为祂圣父王国的伴侣。而[我们的皇帝]常为祂所爱，从上获得帝国权力的渊源，因祂神圣称号的力量而坚强，已在很长一段岁月中统治了世界的帝国。\par

2. 再者，那宇宙的\OriginalTerm{保守者}{Preserver}与祂圣父的旨意一致地安排了天地和天国。同样，祂所爱的我们的皇帝，通过使他在世上统治的人归向那独生圣言和救主，使他们成为祂国度的合适子民。\par

3. 正如那人类共同的\OriginalTerm{救主}{Saviour}，以祂不可见的神圣力量作为\OriginalTerm{善牧}{Good Shepherd}，从祂的羊群那里远远驱走那些曾飞越这地上气域并附着于人类灵魂的\OriginalTerm{背教之灵}{apostate spirits}，如同驱赶野兽。同样，祂的这位朋友，蒙受天上的恩宠胜过所有仇敌，按照战争惯例征服和惩罚真理的公开对手。\par

4. 那位先在的\OriginalTerm{圣言}{Word}，万物的\OriginalTerm{保守者}{Preserver}，将真智慧和救恩的种子分赐给祂的门徒，并立即启迪他们，使他们明白祂圣父国的知识。我们的皇帝，他的朋友，作为天主圣言的解释者，旨在召回整个人类认识天主；向所有人耳中清楚宣告，并以强大的声音向地上所有居民宣布真理和虔诚的律法。\par

5. 再者，普世的\OriginalTerm{救主}{Saviour}为那些从此世走向彼处的人敞开祂圣父国度的天门。我们的皇帝，效法祂的神圣榜样，清除了他尘世统治中一切不敬的谬误污点，邀请每位圣洁虔诚的崇拜者进入他的帝国宫殿，热切渴望拯救那由他担任指定舵手的巨船及其全体船员。而且，在所有执掌罗马帝国权力的人中，唯独他被至高君王赐予三个十年期的统治，现在庆祝这个节庆，并非像他的祖先那样，以敬拜地狱的\OriginalTerm{魔鬼}{demons}或\OriginalTerm{诱惑之灵}{seducing spirits}的显现，或不敬之人的欺诈和欺骗伎俩来庆祝；而是作为对赐予他如此殊荣的祂的感恩，以及承认从祂手中所领受的祝福。他不效仿古代习俗，以血和污秽玷污他的帝国宫殿，也不以火和烟和祭献来讨好地狱的神祇；而是向普世君主献上悦乐而可接纳的祭品，即他自己的皇帝灵魂，以及真正合宜于侍奉天主的心灵。\par

6. 因为唯有这祭献才蒙祂悦纳：我们的皇帝已经学会，以洁净的心思和意念，无需火与血的介入而奉献这祭献，同时他以自己灵魂所蕴藏的真理教义所坚固的虔诚，用华丽的言辞赞美天主，并通过他的帝国行为模仿天主的圣爱\OriginalTerm{神圣的仁爱}{Divine philanthropy}。他完全奉献于天主，将自己作为高贵的祭品、他所受托治理的世界的初果奉献给天主。这第一且最大的祭献，我们的皇帝首先奉献给天主；然后，作为忠信的牧者，他奉献的并非\Quote{著名的头生羔羊的百牲祭}，而是他所照管的羊群的灵魂，即那些理性存在者，他引导他们认识并虔诚敬拜天主。\par

% structure: p0023 source=h2 target=section reason=relative-heading-rank; base=h2

\section{第三章}

1. 祂欣然接纳并欢迎这祭献，并称赞呈献这庄严高贵祭品的人，延长他的统治多年，以更大的恩赐回报皇帝对祂的神圣事奉。因此，祂允许他在帝国普遍繁荣的时期庆祝每一个连续的节日，在每个十年期结束时，让他的一个儿子分享他的帝国权力。\par

2. 长子承袭父亲之名，在第一个十年期结束时成为帝国的共治者；次子在第二个十年期；第三子在第三个十年期，即我们当前节庆的场合。如今第四个时期已经开始，他的统治时间进一步延长，他渴望通过召集更多亲属分享他的权力来扩展他的帝国权威；并通过设立凯撒，实现了圣先知们的预言，正如他们在许多世代以前所说的：\BibleQuote{至高者的圣民必要得国。}\par

3. 这样，全能的上主亲自赐予我们最虔诚的皇帝年岁和子孙的增加，并使他治理万国的权柄依然清新蓬勃，仿佛正在其最初的活力中萌发。正是祂为他指定了这当前的节庆，因为祂使他战胜了所有扰乱和平的敌人；正是祂将他展示为人类真正虔诚的榜样。\par

4. 因此，我们的皇帝如同璀璨的太阳，通过凯撒的临在，仿佛用他自己光芒的远射光线，照亮了他帝国最遥远的臣民。对于居住在东方地区的我们，他赐予了一位与他相称的儿子；第二位和第三位则分别赐予他帝国的其他部分，以成为从他发出的光的明亮反射器。再一次，他仿佛将四位最尊贵的凯撒如同骏马套在同一轭下，置于帝国的战车中，他高高在上，用神圣和谐与一致的缰绳引导他们的方向；而他自身处处临在，洞察每一事件，周游世界的每一区域。\par

5. 最后，他既披戴天上主权的表象，便仰首注视上方，按照那神圣原型的模式塑造他的尘世治理，在效法天主的独一统治中感受力量。而这效法是由普世的上主唯独赐予世上受造物中的人类的：因为唯有祂是主权权力的创造者，祂命定一切都要服从于独一的统治。\par

6. 确实，君主政体远超过其他所有政体和政府形式：因为与之对立的民主平等，毋宁说可被称为无政府状态和混乱。因此，天主是独一的，而非两位、三位或更多：因为断言多神显然就否认了天主的存在。有一位独一的君王；祂的圣言和皇家律法是独一的：这律法不是以音节和词语表达，不是书写或刻在石版上，因而随时间的流逝而朽坏；而是活的、自存的圣言\OriginalTerm{活的、自存的圣言}{living and self-subsisting Word}，祂本身就是天主，并为所有在祂之后并服从祂权柄者管理祂父的王国。\par

7. 祂的侍从是天上的万军：天主千万的使者臣仆；地上以上无法计数的军队；以及在天堂本身内那些无形的灵体，他们的行动用于管理这个世界的秩序。所有这些的统治者和首领是王室的圣言，作为至高君王的摄政。祂被赋予元帅、大司祭、父的先知、大谋略天使、父之光的荣耀、独生子，以及根据圣经作者的神谕\OriginalTerm{圣作者们的预言}{oracles of the sacred writers}所记载的无数其他名号。父既立祂为活的圣言、律法和智慧，一切祝福的圆满，就将这最好最大的恩赐赐予所有服从祂统治的臣民。\par

8. 祂自身渗透万有，处处临在，以毫不吝啬的手向一切揭示祂父的恩赐，甚至将这主权权能的样本赐予地上理性的受造物，因为祂为按自己肖像所造之人的心灵提供了神圣的官能\OriginalTerm{神圣的官能}{Divine faculties}，由此它也能够拥有其他源于同一天上源泉的德行。因为唯有独一的天主是智慧的；唯有祂是本质上良善的；唯有祂是大能的，正义之父\OriginalTerm{正义之父}{Parent of justice}，理性与智慧之父\OriginalTerm{理性与智慧之父}{Father of reason and wisdom}，光明与生命之泉\OriginalTerm{光明与生命之泉}{Fountain of light and life}，真理与德性的施予者\OriginalTerm{真理与德性的施予者}{Dispenser of truth and virtue}；总之，帝国本身以及一切统治权能和权力的创立者\OriginalTerm{帝国的创立者}{Author of empire}。\par

% structure: p0032 source=h2 target=section reason=relative-heading-rank; base=h2

\section{第四章}

1. 然而，人从何处获得这种知识？谁曾将这些真理传达到凡人的耳中？属血肉的舌头从何得能谈论与血肉或物质之性如此迥异的事物？谁曾凝视那不可见的君王，并在他身上见到这些完美？肉体感官可以理解与自己性质相近的元素及其组合，但还没有人夸口能以肉眼鉴察那统御万物的不可见的国度，也没有凡俗之性曾辨别出完美智慧的美丽。谁曾透过血肉的中介看见正义的容颜？人又从何处得来合法主权与帝王权力的概念？一个由血肉组成的存在者，从何产生绝对统治的思想？谁曾向地上的凡人宣告那些不可见、无界定、没有外形的无形实体？\par

2. 诚然，这些事只有一位解释者：那遍及万有的天主圣言。因为祂是存在于人内的理性与理智存在者的创造者；祂自身与祂圣父的天主性合而为一，便将祂圣父恩惠的涌流倾注于祂的儿女。因此，一切人——无论希腊人或蛮族人——都同样拥有自然的、未经教导的思想能力：因此有理性和智慧之觉，正直与正义的种子，生活技艺的理解，德行的知识，智慧的可贵名称，以及对哲学学问的高贵热爱。因此有了一切伟大而美好之物的知识：因此有了对天主自身的理解，以及相称于敬拜祂的生活：因此有了人的王权，以及他对世上受造物不可战胜的主宰。\par

3. 当那作为理性受造者之父的圣言，按照天主的肖像和模样，在人的心灵上印下印记，并使他成为王者的受造物，因为祂唯独赐给他（在所有地上的受造物中）治理和服从的能力（以及在这世上对祂天国所应许之希望的远见和预知，为此祂亲自来临，并作为祂儿女的父亲，不嫌与凡人交谈）；祂继续培育自己所播下的种子，并从以上更新祂恩宠的施予；向众人展示分享祂天国的承诺。因此，祂呼召人，劝勉他们准备好走向天国的旅程，并为自己预备相称于他们召唤的衣服。藉着一种不可言喻的能力，祂的教导充满了世界的每一处，藉着一个尘世的国度的比喻，表达那个祂热切邀请所有人类的天国，并将它呈现为他们值得盼望的对象。\par

% structure: p0036 source=h2 target=section reason=relative-heading-rank; base=h2

\section{第五章}

1. 在这希望中，我们这位蒙天主恩宠的皇帝甚至在此生就已分享，因为他由天主恩赐了本性的德行，并将祂恩惠的涌流收纳于自己的灵魂中。他的理性源自一切理性的伟大源头：他是智慧的、良善的、正义的，因他与完美的智慧、良善、公义相通；他是德行的，因他效法完美德行的典范；他是勇敢的，因他分享了天上的力量。\par

2. 诚然，唯有那按照天上国度的标准，将自己的灵魂塑造为王者德行的人，才配得帝王的称号。但那对这一切福分陌生、否认宇宙的君主、不向众灵的天上圣父效忠的人；那不以相称于帝王的德行披戴自己，却以道德的畸形与卑劣覆盖自己灵魂的人；那以野兽的狂暴取代王者的仁慈，以恶毒之心的不治之毒取代宽厚的性情，以愚妄取代审慎，以鲁莽——所有恶习中最可憎的——取代理性和智慧，因为从中，如同从苦泉，生出最致命的果实，如根深蒂固的放荡生活、贪婪、谋杀、不敬和蔑视天主；一个沉溺于这些恶习的人，无论他因专制的暴力而被视为多么强大，他实在不配称为皇帝。\par

3. 因为，一个人的灵魂刻有上千个假神的荒谬形象，他如何能展现真正天上王权的对应物？或者，一个人使自己臣服于上千个残忍主人的统治，他如何能成为别人的绝对主宰？他是卑下享乐和无法驾驭的欲望的奴隶，是不义榨取的财富的奴隶，是愤怒和情欲的奴隶，也是怯懦和恐惧的奴隶；是无情魔鬼和毁灭灵魂的邪灵的奴隶。\par

4. 那么，让我们的皇帝，凭着真理本身的见证，被宣布为唯一配得这称号的人；他亲切于至高君主自己；唯有他是自由的，不，他真是主宰：高于财富的渴望，超于性欲；甚至战胜了本性的快乐；控制而非被愤怒和情欲所控制。他确实是皇帝，拥有与他的事迹相称的称号；是真理中的胜利者，他战胜了那些征服其余之人的情欲；他的品格是按照至高君主的神圣原型形成的，他的心灵如同镜子般反映着德行的光辉。因此，我们的皇帝在明辨、良善、正义、勇敢、虔诚、对天主的敬爱上是完美的；他真正且唯独是哲学家，因为他认识自己，并完全意识到每一种福分的供应都是从他自己之外的一个源头——甚至是从天上——倾注于他的。他藉着袍服的华彩宣告至高权柄的威严称号，只有他值得穿戴那如此合身的帝王紫袍。\par

5. 他实在是一位皇帝，日夜以祈祷呼求天上圣父的眷顾，他的热切渴望专注于天上的国度。因为他知道，现今的事物，如同河流一般流逝消逝，易于朽坏和死亡，不值得与万有的君宰相比；因此他渴望天主那不朽、无形的王国。他相信必得这王国，因为他的思想已提升到穹苍之上，心中充满无法言喻的渴望，向往那里闪耀的荣光，与此相比，他视今世的珍宝为黑暗。因为他看出地上的主权不过是短暂、易逝的统治，掌握着必死且暂时的人生，他认为这并不比牧羊人、牧牛人或牧人的权柄高多少：甚至比他们的更沉重，因为要统治更倔强的臣民。人民的欢呼和谄媚的声音，他视为麻烦多于喜悦，因为他性格坚定，心智受过真正的训练。\par

6. 再次，当他看到臣民的军役，他的大军列阵，无数的骑兵和步兵都完全服从他的命令时，他不感到惊奇，也不因拥有如此强大的力量而骄傲；而是转向内省，认识到同样的共同本性。他微笑地看着自己绣金和饰花的衣袍，以及皇帝的紫袍和冠冕，当他看到众人如同孩子见到怪物一样，惊奇地凝视这华丽的景象时。他超越这些情感，用认识天主的衣袍装饰自己的灵魂，这衣袍的绣花是节制、正义、虔诚和所有其他德行；这样的衣袍才真正适合一位君主。\par

7. 别人如此渴望的财富，如黄金、白银或珍宝，他视之为本身不过石头和无价值的物质，不能保护或防御邪恶。因为这些事物有什么力量能使人免于疾病，或阻止死亡的临近呢？他通过亲身使用这些事物而知道这一真理，因此他漠然地看待臣民华丽的服饰，并嘲笑那些被这些事物吸引的人的幼稚。最后，他戒除一切饮食的过度，把多余的珍馐留给贪食者，认为这样的放纵尽管对别人合适，对他却不合适，并且深信它们有害的倾向，以及对灵魂理智能力的蒙蔽效果。\par

8. 由于这一切原因，我们这位受天主教导、思想高尚的皇帝，追求今生之上更高的目标，呼求天上的圣父，如同渴望他的国度；在他生命中的每一个行动都表现出虔诚的精神；最终，如同一位智慧而良善的导师，将认识万有的至高主宰传授给他的臣民。\par

% structure: p0045 source=h2 target=section reason=relative-heading-rank; base=h2

\section{第六章}

1. 而天主自己，作为未来赏报的保证，现在赐给他如同由繁荣的时期组成的三十年冠冕；如今，在三个十年的循环之后，他准许全人类庆祝这个普世的节庆，更确切地说，这个普世节庆。\par

2. 当地上的人如此喜乐，仿佛被戴上认识天主的鲜花冠冕时，我们或许有理由认为，天上的歌咏团也因自然的同情而与人间的喜乐相结合；不，至高的君主自己，如同慈爱的父亲，因孝顺子女的敬拜而喜悦，并因此乐意延长他们服从的创始者和原因的统治时期；他不仅将他的统治限于三个十年的循环，而且将其延伸至最遥远的时期，甚至到遥远的水恒。\par

3. 永恒在其整个范围内是超越朽坏和死亡的：它的起始和范围都不能被人的思想所测度。它不允许其中心点被感知，也不允许所谓当前持续的时间被探究的心灵所把握。那么未来和过去更是如此：因为过去已不存在，已经逝去；而未来尚未到来，因此也不存在。至于所谓现在的时间，它甚至在我们思考或说话时就已经消失，比话本身说得还快。因此不可能以任何方式把握这个时间为现在；因为我们要么期待未来，要么默想过去；现在从我们滑过，在我们思考的瞬间就逝去了。因此，永恒在其整个范围内，拒绝并拒不服从人的理性。\par

4. 但它却不拒绝承认自己的君宰和主，并承载他，如同他骑在它上面，以他赋予的美丽装饰为乐。而他自己，并非像诗人想象的那样用金链捆绑它，而是用不可言喻之智慧的缰绳控制它的运动，和谐地调整了它的月、季、时、年，以及昼夜的交替，从而给它附上了各种界限和度量。因为永恒，按其本性是直的，延伸至无限，并因其恒久的存在而得名为永恒，在其所有部分中相似，或者更确切地说，没有划分或距离，只沿着直线延伸。但是天主，通过中间的段落将其分布，并将其像一条很长的线一样在许多点上划分，在其中包含了大量的部分；尽管它本性是一，并像统一本身，他却给它附上了众多数字，并赋予它本身无形却无尽的形式。\par

5. 因为他首先在其中塑造了\OriginalTerm{无形质料}{formless matter}，作为一种能够接受一切形式的实体。接着，他藉着二的力量，赋予质料以性质，并使那先前毫无\OriginalTerm{优美}{grace}的事物具有了美。再者，藉着三，他塑造了一个由质料和形式组合而成的身体，呈现长、宽、深三维。然后，从二的加倍，他设计了\OriginalTerm{四元素组}{quaternion}，即土、水、气、火，并命它们成为供应这宇宙的永恒之源。再者，四产生十。因为一、二、三、四的总和是十。三乘以十得出一月的周期；十二个连续的月完成太阳的行程。由此，年岁的循环和季节的变迁，如同绘画中的色彩变化，赋予那先前无形且缺乏美的永恒以优美，为那些命定在其中穿行生命历程者带来舒缓和喜悦。\par

6. 正如赛场因规定的距离而为那怀着得奖希望奔跑的人划定界限；又如远行旅客的道路以驿站和分隔的里程为标志，以免旅行者因无尽的远景而气馁；同样，宇宙的主宰以自身智慧的约束力掌控永恒本身，按他认为最好的方式引导并转变其进程。我说，正是同一位天主，如此将先前\OriginalTerm{无定限的永恒}{undefined eternity}装扮，如同用美丽的色彩和盛开的鲜花，以阳光照亮白昼；并当夜幕笼罩黑暗时，却使闪耀的星辰如金片在其中发光。是他点明了晨星的璀璨、月亮的变幻光华以及灿烂的星群队伍，并将穹苍的展开如同巨大的披风，以各种美丽的色彩装饰。再者，他创造了高深广阔的空气，使世界在长宽上感受到它的清凉影响，并命空气本身以各种飞鸟为装饰，留下这广阔的空间之洋供一切可见或不可见的、穿行于天穹之域的受造物翱翔。在这大气之中，他使大地悬于其中心，并以海洋环绕，如同美丽的天蓝色外衣。\par

7. 他命定这大地同时成为其中所有受造物的家园、养育者和母亲，并以雨和泉水浇灌它，使它滋生各种植物和花朵，以供生命的享受。当他按自己的肖像造了人——地上受造物中最尊贵、最亲近于他的，一个赋有理智和知识的受造物，理性与智慧之子——他将统治权赐予人，管辖一切在地上活动生存的其他动物。因为人实在是所有地上受造物中最蒙天主喜爱的：人，我说，慈父般的天主将兽类归服于他；为了他，天主使海洋可通航，并用各种植物为王冠装饰大地；赐予他推理能力以获得一切知识；甚至将深渊中的生物和空中有翼的居民置于他的控制之下；允许他默观天上之物，并启示太阳和月亮的运行与变化，以及行星和恒星的周期。简言之，在地上众生中，唯独人领受了诫命，要承认天主为他的天上之父，并赞美他为永恒自身的至高主宰。\par

8. 但造物主以一年四季限制了\OriginalTerm{不变的永恒进程}{unchangeable course of eternity}，以春天的临近终结冬天，并以均等的平衡调整那开始年度周期的季节。他既以春天的各种产物美化了永恒的时间进程，又添加了夏日的炎热；然后以秋天的间隔赐予劳作的休息；最后，以冬日的阵雨清新并洁净季节，使其经充沛雨水后变得光滑亮泽，如同高贵的骏马，再次带到春天的大门前。\par

9. 至高主宰既以这些智慧的绳索将自身的永恒与年度的循环如此联结，便将其托付给一位伟大的治理者，即他的\OriginalTerm{独生圣言}{only begotten Word}，他作为\OriginalTerm{万有的保存者}{Preserver of all creation}，获得宇宙大权的缰绳。圣言从仁慈的父领受这产业，将穹苍上下的一切在和谐的整体中结合，引导它们一致的运行；以完全的正义提供一切有利于地上理性受造物的事物，为人的生命划定指定的限度，并同样准许众人甚至在此时预先开始未来的存在。因为他已教训他们：在这现世之外，有一种神圣而蒙福的存在状态，为那些在此世以天上福佑之希望得支持者保留；并且那些度了有德和虔诚生活的人将由此移居到更美好的居所；同时，他判定那些在此世有罪和邪恶的人按其罪行受惩罚的地方。\par

10. 再者，如同在公共赛会中分发奖品，他向胜利者宣告各种冠冕，并以不同德行的赏报装饰每人；但对于我们善良的皇帝，他穿着虔诚的外衣，他宣称一份更高的酬报已为他劳苦预备；作为这酬报的前奏，他准许我们现在聚集在这个由\OriginalTerm{完全数}{perfect numbers}——\OriginalTerm{三倍的十}{decades thrice}，以及\OriginalTerm{十倍的三个一组}{triads ten times repeated}——构成的节日。\par

11. 这些数中的第一个，即\OriginalTerm{三元}{triad}，是\OriginalTerm{一}{unit}的产物，而一本身就是数之母，并掌管所有月份、季节、年份及一切时间段。的确，它有理由被称为一切数的起源、基础和原理，其名称源于它不变的特质。因为，当其他每个数因增减而变化时，唯独一保持固定和稳定，远离一切众多及其所生成的各种数，如同那不可分割的本质，它与所有其他事物相异，但万物凭借参与它而得以存在。\par

12. 因为一是一切数的创始者，一切的众多都是由一的组合和相加而成；没有一，根本不可能设想数的存在。但一本身独立于众多，超越并高于一切数；它确实形成并制造一切，却不从任何事物接受增加。\par

13. 与一相近的是三元；同样不可分割且完美，是由偶数和奇数构成的第一种和数。因为完美的数二，加上一，便形成三元，即第一个完美的合成数。三元通过解释何为平等，首先教导人们正义，它本身具有相等的起点、中点和终点。它也是神秘、至圣而君尊的\OriginalTerm{天主圣三}{Trinity}的肖像，天主圣三虽无始无终，却包含一切受造物存在的萌芽、理性和原因。\par

14. 因此，三元的力量有理由被视为万物的第一因。再者，数字十包含一切数的终结，并将它们终止于自身，确实可被称为丰满而完美的数，因为它包含一切种类的数和一切度量、比例、和谐与协和。例如，一通过相加形成并在数字十处终止；它们以十为母，又如同其行程的界限，将十视为其历程的终点。\par

15. 然后它们进行第二次循环，又是一次第三次、第四次，直到第十次，如此通过十个十进组构成第一百个数。从那里返回最初的起点，它们再次进到数字十，当十次完成第一百个数后，它们再次后退，围绕同一障碍进行漫长的行程，以循环运动从自身回到自身。\par

16. 因为一是十的十分之一，十个一组成一个十进组，后者本身就是一的界限、确定的目标和边界：它终止数的无限性；是诸一的终点和结束。再者，三元与十进组结合，并进行三次十的循环，便产生最自然的数三十。因为三元相对于一，正如三十相对于十。\par

17. 它也是那仅次于太阳光辉之天体的恒定行程界限。因为月亮从一次与太阳会合到下一次会合的行程，完成一个月的周期；此后，它仿佛获得新生，重新开始新的光芒和新的日子，以三十个一、三个十进组和十个三元组被装饰和尊崇。\par

18. 同样，我们得胜的皇帝的普世统治被万善之赐予者所尊荣，如今进入一个新的福佑领域，目前完成这\OriginalTerm{三十年庆典}{tricennalian festival}，但由此向前展望更遥远的时段，并怀有对未来在天国中福佑的希望；在那里，不是单一的太阳，而是无数光明（众）围绕全能之主，每一位都超越太阳的光辉，因来自永恒光源的光线而荣耀灿烂。\par

19. 在那里，灵魂享受着美好而不朽的福佑，享受着一种超越忧愁的生命；在那里，有纯洁而神圣的快乐体验，以及无法度量、无尽延伸的时间，延至无垠的空间；不按日月的间隔、年岁的循环或时节的重现来界定，而是与一个无终的生命相称。这生命不需要太阳之光，也不需要月亮或星辰的光辉，因为它拥有那伟大的光明者本身，即天主\OriginalTerm{圣言}{Word}，全能之主的独生子。\par

20. 因此，神秘而神圣的谕示启示他为\OriginalTerm{义德的太阳}{Sun of righteousness}，以及远超一切光的\OriginalTerm{光}{Light}。我们相信，他也用正义的光线和智慧的光辉照亮那三重福佑的天上诸能，并且他接纳真正虔诚的灵魂，不仅在天堂的境界内，而且进入他自己的怀抱，并确证他自己所赐予的应许。\par

21. 没有肉眼看见过，也没有耳朵听见过，人心在血肉之躯里也不能明了天主为那些在此世以虔敬\OriginalTerm{恩宠}{graces}所装饰的人所预备的是什么；这些福佑也等待着您，至虔敬的皇帝，因为自世界创立以来，惟有您，宇宙的全能之主授予了净化人类生活的权柄：并且祂向您启示了祂自己的\OriginalTerm{救恩标记}{symbol of salvation}，藉此祂克服了死亡的力量，战胜了一切仇敌。而这胜利的\OriginalTerm{纪念标志}{victorious trophy}，邪灵的鞭笞，您用来对抗偶像崇拜的谬误，不仅战胜了所有不虔敬的蛮族仇敌，也战胜了同样野蛮的对手——邪灵本身。\par

% structure: p0067 source=h2 target=section reason=relative-heading-rank; base=h2

\section{第七章}

1. 因为我们由两种不同的本性构成，我指的是身体和精神，其中一种对所有人可见，另一种不可见，针对这两种本性，有两种野蛮而凶残的仇敌不断列阵，一种以不可见的方式，另一种以公开的方式。前者以身体的力量对抗我们的身体，后者以无形的攻击围攻赤裸的灵魂本身。\par

再者，可见的野蛮人，如同荒野的游牧部落，与野兽无异，袭击文明民族的疆域，蹂躏其国土，奴役其城市，如无情的沙漠之狼般扑向居民，摧毁一切落入其权势之下者。但那些更甚于野蛮人的无形仇敌——即那些穿行于大气区域的、毁灭灵魂的\OriginalTerm{魔鬼}{demons}——已通过卑劣的\OriginalTerm{多神论}{polytheism}的陷阱，成功地奴役了全人类，致使他们不再认识真天主，反而迷失于无神论的错误迷宫之中。因为他们从某处——我不知从何处——弄来了从未存在过的神祇，却将那唯一而真实的天主置于一旁，仿佛他不存在一般。\par

因此，肉身的生殖被他们奉为神祇，其对立面——即肉身的解体与毁灭——同样被神化。前者作为生殖力的源泉，以\OriginalTerm{维纳斯}{Venus}之名受到礼敬；后者因富有并对人类拥有强大统治权，被称为\OriginalTerm{普路托}{Pluto}和\OriginalTerm{死亡}{Death}。因为那个时代的人，除了自然产生的生命之外一无所知，便宣称这生命的因由和起源是神圣的；又因不相信死后有任何存在，便宣告死亡本身是全能的征服者和伟大的神。因此，由于不知有责任——仿佛注定要被死亡消灭——他们过着名不副实的生活，行径当受千次死亡。他们心中没有天主的念头，没有神圣审判的期待，没有对灵魂存在的记忆或反思；他们承认一个可怖的至尊——死亡，确信身体被其力量解体便是最终的消亡，便赐予死亡\Quote{伟大而富有之神}的称号，故有\OriginalTerm{普路托}{Pluto}之名。如此，死亡便成了他们的神；不仅如此，凡是他们认为与死亡相比珍贵、且有助于生活享乐之物，亦皆如此。\par

于是，肉欲之乐成了他们的神；滋养及其之物成了神；树木的果实成了神；醉酒狂欢成了神；肉欲的渴望与快乐成了神。故有\OriginalTerm{色列斯}{Ceres}与\OriginalTerm{普洛塞庇娜}{Proserpine}的秘仪，后者被\OriginalTerm{普路托}{Pluto}掳走及随后归还；故有\OriginalTerm{巴克斯}{Bacchus}的狂欢，以及\OriginalTerm{赫拉克勒斯}{Hercules}被醉酒击败，如同被更强的神征服；故有\OriginalTerm{丘比特}{Cupid}与\OriginalTerm{维纳斯}{Venus}的淫逸仪式；故有\OriginalTerm{朱庇特}{Jupiter}本人迷恋妇女与\OriginalTerm{伽倪墨得斯}{Ganymede}；故有关于沉沦于女色与享乐的诸神的放荡传说。\par

这便是迷信的武器，这些残暴的野蛮人和至高天主的仇敌用以折磨——甚至完全征服——人类的武器；他们在各处竖立不敬的纪念碑，在每个角落建造虚假宗教的神龛和庙宇。\par

非但如此，当时的统治势力竟被错误的力量如此奴役，以致用自己同胞和亲族的血来安抚他们的神祇；竟磨刀霍霍，挥向那些挺身捍卫真理的人；竟发动无情的战争，举起不洁之手——不是对外邦或蛮族的敌人，而是对与他们有家庭和亲情纽带的人，对弟兄、亲族和最亲爱的朋友——这些朋友决心在德行和真虔诚中，尊崇并敬拜天主。\par

这些君王便是以如此疯狂的精神，将那些奉献给万王之王的人祭献给他们的魔鬼神祇。另一方面，他们的牺牲者，作为真虔敬事业的高贵殉道者，决心欣然迎接光荣的死亡胜过生命本身，全然蔑视这些残暴。作为天主的士兵，他们以忍耐的力量刚强，讥笑各种形式的死亡：火、剑、十字架的酷刑；暴露给野兽、沉入深海；肢体的斩断与烧灼、挖眼、全身的残害；最后是饥饿、矿山的苦役和囚禁——因着他们对天上君王的热爱，所有这些苦难他们都认为胜过任何世俗的美好或快乐。同样，妇女也表现出不逊于男子的坚定与勇敢精神。\par

一些妇女与男子同受这些争战，并获得相同的德行报偿；另一些妇女被强行掳去，成为暴力与污辱的牺牲品，却宁死也不愿受辱；还有许多妇女，甚至未曾忍受省督用以恐吓她们的威胁，就勇敢地承受了各种酷刑和各式死刑。如此，全能君王的这些英勇士兵，以灵魂的坚定刚毅，对抗多神论的敌对势力；如此，这些天主的仇敌、人类救恩的反对者——比凶猛的野蛮人更为残忍——以人血为祭品而自娱；如此，他们的仆役仿佛饮尽了为所侍奉的魔鬼而行的不义杀戮之杯，为他们预备了这可怕而不虔的筵席，使人类走向毁灭。\par

在此悲惨境况下，这些受苦者的天主和君王应当如何行事？他岂能漠不关心他最亲密朋友的安全，或在这极大危难中抛弃他的仆人？诚然，无人会认为那不去努力拯救同船水手、任凭船只与全体船员沉没的人是一位谨慎的舵手；也绝无将军会如此鲁莽，不加抵抗就将自己的盟友交给敌人任意处置；忠实的牧羊人也绝不会漠视羊群中一只羊的迷失，反而会将其余的羊留在安全之处，为那迷失的羊冒险行事，甚至必要时与野兽搏斗。\par

10. 然而，万有的伟大主宰的热忱并非为了无灵的羊群：他的眷顾施予他忠信的军旅，那些为他而奋战的人。他亲自嘉许他们为敬虔事业所做的争战，并尊荣那些回到他面前的人，赐予唯独他能赐予的胜利奖赏，将他们联合于天使的唱诗班。其他人他仍保存在世上，将敬虔的活种传给后代；使他们既成为他报复不敬虔者的目击者，也成为这些事件的讲述者。\par

11. 此后，他伸臂审判仇敌，以神圣义怒的打击彻底毁灭了他们，迫使他们无论多么不情愿，也要亲口承认并弃绝自己的邪恶；却将从地上高举、荣耀地抬举那些长久受压迫、被众人弃绝的人。\par

12. 至高主宰的作为正是如此，他命定一位无敌的战士作他天赐报复的执行者（因我们皇帝的卓越虔诚喜爱\Quote{天主的仆人}这一称号），并使他战胜一切反对者，兴起他一人对抗众多仇敌。因为仇敌确实无数，是许多恶灵的朋友（虽然实际上他们毫无所有，因此如今不复存在）；但我们的皇帝是独一的一位，由全能独一主宰所任命，是他的代表。而他们，以不虔的精神，用残忍屠杀摧毁义人；但他，效法他的救主，只知拯救人的性命，竟宽恕了不虔者本人并教导他们敬虔。\par

13. 于是，他真正配得\Quote{胜利者}之名，征服了两类野蛮人；以明智的使节安抚野蛮部落，迫使他们认识和承认他们的尊长，将他们从无法无天的野蛮生活引向理性与人性的管治；同时他以事实证明了凶猛无情的无形灵类早已被更高能力击败。因为宇宙的守护者以无形的审判惩罚了这些无形之灵；而我们的皇帝，作为至高主宰的代表，继续这场胜利，掠走那些早已死亡化为尘土者的战利品，并将掠物慷慨分给他胜利之主的士兵。\par

% structure: p0081 source=h2 target=section reason=relative-heading-rank; base=h2

\section{第八章}

1. 因为他一明白无知的群众被这些用金银制成的错误妖怪激起虚妄幼稚的恐惧，他就认为应当移除这些，如同放在黑暗中行走之人路上的绊脚石，并从此为所有人开辟一条平坦无障碍的御道。\par

2. 作出这个决定后，他认为镇压这邪恶不需要任何士兵或军事力量：他少数几个朋友就足够这项任务，他仅以简单的意志表达派遣他们巡视各个省份。\par

3. 于是，他们因对皇帝虔诚和自身对天主的虔敬而满怀信心，穿过无数部落和民族，在每个城市和乡村废除这古老的错误体系。他们命令司祭们自己在普遍的嘲笑和轻蔑中，将他们的神祇从黑暗隐蔽处带到光天化日之下。然后他们剥去神像的装饰，向众人展示隐藏在外表彩绘下丑陋的真相；最后，他们刮下任何看起来有价值的材料，用火熔化以检验其价值，之后将他们认为合用的部分保管并分开，将完全无用的部分留给迷信的崇拜者，作为他们羞耻的纪念物。\par

4. 同时，我们可敬的皇帝本人也投身于类似我们描述的工作。因为就在这些昂贵的死者偶像被剥去珍贵材料的同时，他也攻击了那些黄铜制的偶像；下令用绳索将它们从原处拖走，仿佛俘虏般带走，那些神话的愚昧曾将它们尊为神祇。我们威严的皇帝接下来的关怀是仿佛点燃一支明亮的火炬，藉其光芒他环视帝国的目光，查看是否还有隐藏的谬误痕迹。\par

5. 如同目光锐利的鹰向天飞翔时，能从其高处看见地上最遥远的事物：同样，当他驻跸于自己美丽城市的皇宫时，竟如从瞭望塔发现\Place{腓尼基}省一个隐藏致命的灵魂陷阱。那是一个树林和庙宇，并非位于任何城市中央或公共场所（如为壮观效果通常如此），而是在偏僻少人行走的道路旁，在\Place{黎巴嫩山}的部分山顶，献给了名为维纳斯的污秽恶魔。\par

6. 那是一个邪恶学府，为所有放纵于不洁和以柔靡毁坏身体者而设。在此，不配称男的人忘记了自己性别的尊严，以柔靡行为讨好恶魔：在此也有妇女的非法交易和奸淫，以及其他可怖无耻的行径，在这庙宇中犯下，如同在一个超越法律范围和约束的地方。\par

同时，这些邪恶因没有观察者而未被制止，因为品性端正者无人敢涉足此类场所。\par

7. 然而，这些行为未能逃脱我们威严皇帝的警惕，他亲自以特有的先见审视它们，认为这样的庙宇不宜接受天空之光，下令将建筑物及其供物彻底摧毁。于是，遵从皇帝的法令，这些不洁迷信的器械被立即废除，军事力量被用来净化这地方。如今，那些从前生活毫无节制的人，因皇帝惩罚的威胁，学会了实行自制。\par

8. 如此，我们的皇帝将这虚幻邪恶体系的伪装撕下，公之于众，同时向所有人公开宣告他救主的圣名。没有辩护者出现；被揭露的骗局作者们既得不到神祇或邪魔、先知或占卜者的帮助。因为人的灵魂不再被浓重的黑暗笼罩：而是被真虔诚的光芒照亮，他们为祖先的愚昧而哀叹，为祖先的盲目而惋惜，同时为自己从如此致命的错误中得救而欣喜。\par

9. 如此，按照全能天主的旨意，并通过我们皇帝的中介，每一个敌人，无论是可见的还是不可见的，都被彻底驱除；从此和平——青春的幸福养育者——将她的统治延伸到全世界。战争不再发生，因为诸神不再；不再有田野或城镇的战争，不再有人血的流淌，如同从前邪魔崇拜和偶像崇拜的疯狂盛行时那样折磨人类。\par

% structure: p0092 source=h2 target=section reason=relative-heading-rank; base=h2

\section{第九章}

1. 如今我们可以将现在与过去相比，回顾这些幸福的变化，与过去的邪恶形成对比，并注意在古代，人们为这些假神精心准备了门廊、圣地、树林和庙宇，并用丰盛的供品装饰他们的神殿。\par

2. 当时的最高统治者确实对诸神的崇拜非常重视。属下的国家和人民也根据祖先的宗教习俗，在乡村和每个城市，甚至在他们的房屋和密室中，用雕像尊敬它们。然而，这种虔诚的果实，远非我们今天看到的和平和谐，而是以战争、战斗和叛乱的形式出现，这些贯穿他们一生，并用血和国内屠杀淹没了他们的国家。\par

3. 再者，他们崇拜的对象能以狡猾的谄媚向这些统治者承诺预言、神谕和对未来的知识：然而他们却不能预言自己的毁灭，也不能预知即将到来的覆灭：这无疑是他们欺骗的最大和最令人信服的证明。\par

4. 那些曾经其话语令人敬畏和惊叹的人中，没有一个宣布了人类救主的荣耀来临，或他带来的神圣知识的新启示。无论是\OriginalTerm{皮提乌斯}{Pythius}本人，还是那些强大的神祇，都不能预知他们即将来临的荒凉；他们的神谕也不能指向那个将是他们的征服者和毁灭者的人。\par

5. 哪一个先知或占卜者能预言，他们的仪式会在一个新神祇出现在世界时消失，并且全能君主的认识和崇拜会自由地赐给全人类？他们中谁预知了我们得胜皇帝的威严而虔诚的统治，或他到处战胜假邪魔的胜利，或他们高处的倾覆？\par

6. 哪一个英雄曾宣布，无生命的雕像从无用的形式熔化和转化为人类必要用途？哪一个神祇还有能力说出他们自己的雕像这样被熔化并轻蔑地打成碎片？\par

7. 保护的力量在哪里，以致它们不介入拯救它们被人类摧毁的神圣纪念物？我问，那些曾经维持战争纷争，现在却看着他们的征服者安然处于最深和平中的人在哪里？那些在盲目愚昧的自信中坚持自己，并把这些虚妄当作神祇来信赖的人在哪里？然而，在他们迷信错误最盛之时，并与真理的捍卫者维持着无情的战争时，他们却以与其罪行相称的命运灭亡。\par

8. 那手臂转向天本身的巨人种族在哪里？那些舌头用不敬的话指向全能君主的蛇的嘶嘶声在哪里？这些万主之主的敌对者，依靠众多神祇的帮助，以强大的军力进攻，前面有一些死者的形象和无生命的雕像作为防御。另一边，我们的皇帝，以虔诚的铠甲为保障，以救恩和赋予生命的标记对抗敌人的数量，同时作为对敌人的威慑和对一切伤害的保护；并同时战胜了敌人和他们所侍奉的邪魔，归于他胜利的作者，以感恩和赞美的感恩精神，用纪念碑和言语向全人类宣告这个胜利的标记，将它作为对抗每一个敌人的强大胜利纪念碑立在帝都之中，并明确命令所有人承认这个永不朽坏的救恩象征是罗马权力和世界帝国的保障。\par

9. 这是他普遍给予臣民的训示；尤其是给他的士兵，他劝告他们不要信赖他们的武器、盔甲或体力，而要承认至高天主是每一种善行和胜利本身的赐予者。\par

10. 就这样，皇帝本人（这事实似乎奇怪而难以置信）成为了他军队在宗教操练中的导师，教导他们按照神圣法令献上虔诚的祈祷，向上高举双手，更高举他们心灵的目光到天上的君王，他们应呼求他作为胜利的作者、他们的保护者、守护者和帮助者。他还命令应将一天视为宗教崇拜的特殊机会；我指的是那真正是万日之首和首要的日子，我们主和救主的日子；那日子的名字与光、生命、不朽和一切善美相连。\par

11. 他为自己规定了同样的虔诚行为，在皇宫的内室中敬礼他的救主，遵照神圣的命令履行他的敬礼，并通过聆听圣言来充实他的心智。他整个家庭的管理都托付给了那些献身于天主、以严肃的生活和其他一切美德著称的仆役；而他忠实的卫队，对他的感情和忠诚坚定不移，在他们皇帝身上找到了敬虔生活的导师。\par

12. 再者，他对那胜利记号的尊崇是基于他对其神圣效能的实际体验。在此之前，他敌人的军队消失了；藉此，无形邪灵的力量被驱散；通过它，天主仇敌的狂傲夸耀归于无有，亵渎和咒骂之人的舌头被堵住。藉此记号，蛮族部落被征服；通过它，迷信欺诈的仪式得到了公正的斥责；藉此，我们的皇帝，如同偿还神圣的债务，通过在世界各地竖立其价值的胜利纪念碑，建造了皇家规模的圣殿和教堂，并命令所有人共同建造神圣的祈祷之所，成就了所有善事中的至善。\par

13. 因此，我们皇帝慷慨的这些显著凭证立即出现在帝国的行省和城市中，并很快在每一个国家中闪耀着光辉；这些令人信服的纪念物，斥责并推翻了那些不久以前还疯狂地胆敢与天主作战的不敬神的暴君，他们像狂暴的恶狗一样，将无法对准天主的愤怒发泄在无意识的建筑物上，推倒了祈祷之所，甚至捣毁了其根基，使它们呈现出一座被占领并被遗弃给敌人的城市的景象。这就是那邪恶精神的展示，他们藉此仿佛要攻击天主本身，但很快便体验到他们自己的疯狂和愚昧的结果。没过多久，天堂不悦的风暴一阵吹过，就将他们完全扫除，没有留下亲族、后裔，或他们在人间的任何纪念：因为所有这些人，无论多少，都在一瞬间因神性的报应而消失了。\par

14. 这些狂怒的天主仇敌的命运便是如此；但那位以救恩的十字架武装起来、独自对抗他们的人（不，并非独自，而是得到那唯一主宰的同在和能力的帮助），将毁坏的建筑以更大的规模重建，使后者远远超过前者。例如，除了在以其名字命名的城中建造各种圣殿以荣耀天主外，还在比提尼亚的首都建造了另一座最大最辉煌的圣殿，并在其他各行省的主要城市中以其类似的结构而著称。\par

15. 最重要的是，他在帝国东部选择了两个地方：一个在巴勒斯坦（因为生命之泉从那里流出，如同从源泉中流出，造福万民），另一个在东部的那座以其名源自安提约古的大都市；在那里，作为那部分帝国的首府，他奉献了一座规模无与伦比、美丽绝伦的圣殿给天主。整座建筑被一个广大的围墙环绕，圣殿本身高高耸立，呈八角形，四周环绕着许多厅堂和庭院，并以最丰富的装饰点缀。\par

16. 这就是他在这里的工程。再者，在巴勒斯坦行省，在那曾经是希伯来君主国所在地的城市中，就在主坟墓的原址上，他建造了一座宏伟的圣殿，并以富丽堂皇的宏伟装饰了一座奉献给救恩的十字架的圣所，尊崇那永恒的纪念物，以及救主战胜死亡之权的十字架，以语言无法形容的荣光。\par

17. 在同一国家，他发现了三个因三座神圣洞穴的地点而可敬的地方：他也用昂贵的建筑装饰了它们，对第一次显现救主临在的场景致以恰当的敬意；而在第二座洞穴，他纪念了救主从山顶的最后升天；并在第三座洞穴中，庆祝了他的大搏斗以及加冕的胜利。这一切场所，我们的皇帝如此装饰，是希望向全人类宣扬救赎的象征：\par

18. 那十字架确实回报了他的虔诚热忱；通过它，他的家室和宝座都得以兴旺，他的统治在漫长的岁月中得到巩固，美德的赏报被赐予了他的高贵儿子、亲属以及他们的后裔。\par

19. 确实，这是他侍奉的天主大能的明证，他以公正的手持握天平，分配给每个人应得的赏报。至于那些祈祷之所的破坏者，他们不敬行为的惩罚紧随着他们：他们立刻被扫除，没有留下族裔、家室或家庭。另一方面，他，即在其每一行动中都彰显对主之虔诚敬爱，为他建造皇家圣殿、通过神圣奉献向天下臣民宣扬他名的人——我说，他配得地体验到天主是他皇族家室的保存者和保护者。如此，天主的作为已清楚彰显，而这正是通过那救恩记号的神圣效能。\par

% structure: p0112 source=h2 target=section reason=relative-heading-rank; base=h2

\section{第十章}

1. 对于那些通晓我们神圣宗教奥迹的人而言，关于这救恩的记号，实在有许多可说的。因为它实在是救恩的象征，说得神奇，想得更神奇；它在地上的出现，从一开始就将所有虚假宗教的虚构投入最深的阴影，将迷信的错误埋葬在黑暗和遗忘中，并向所有人启示了那光照人类灵魂的神性之光，即对唯一真天主的认识。\par

2. 因此，普世的转向改善，引导人藐视无生命的偶像，践踏属魔鬼的非法礼仪，并嘲笑先祖们历来的愚行。因此，在各地建立了神圣学问的学校，人在其中受教于救恩真理的规诫，不再惧怕肉眼所见之物，也不以惊奇的目光注视太阳、月亮或星辰，而是承认万有之上的那一位，即那不可见、创造万有的存在，并学习唯独朝拜祂。\par

3. 这就是这伟大而奇妙的标记给人类带来的福分；凭此标记，昔日的邪恶如今不复存在，而前所未闻的德行，以真正敬虔的光辉，处处照耀。\par

4. 劝勉、训诲和鼓励过虔诚圣洁生活的言论，在万民耳中宣讲。不但如此，皇帝本人也亲自宣讲：这位伟大的君王，如同全能君王的代言人，在全世界面前扬声，邀请各国臣民认识真天主，这实在是一件奇事。\par

5. 昔日不虔者妄语充斥皇宫，如今已不复闻；司铎和虔诚的天主崇拜者一同以君王般的赞美诗颂扬祂的尊威。宇宙独一至高统治者的名号传扬于万民；喜讯的福音将人类与全能之王联结，宣告天上的父对其地上子女的恩宠与慈爱。\par

6. 祂的赞歌处处凯旋奏响；凡人的声音与天上天使歌咏团的和谐融为一体；理性的灵魂以身体为乐器，奏出合宜的赞美和朝拜之音献给祂的名。东西方的万民同时领受祂的训诲；南北方的众民在相同原则和律法的影响下，同心合意追求敬虔的生活，赞美独一至高天主，承认祂的独生子、我们的救主为一切福泽的源头，并承认我们的皇帝为地上唯一的统治者，与他虔诚的儿子们一同治理。\par

7. 他本人如同一位技艺高超的舵手，高居国家之舵，以无误的航向引领船只，如同顺风般将百姓带入安全宁静的港湾。同时，天主自身，大君王，从天上伸出祂大能的右手予以保护，赐他战胜一切仇敌，以长久年岁巩固他的帝国；祂还将赐予他更高福分，并证实祂所有应许的真实。然而这些事我们此刻不能细述，当等候进入更美世界的转变：因为属天主的事，不是血肉之眼和耳朵所能完全领悟的。\par

% structure: p0120 source=h2 target=section reason=relative-heading-rank; base=h2

\section{第十一章}

1. 得胜而强有力的君士坦丁啊，在这篇以全能君王的荣耀为崇高主题的讲辞中，请容我向你陈述一些神圣真理的奥秘：并非自以为能教导你这得天主教训的人，也非向你揭示那些他并非藉人、而是藉我们共同的救主，并藉他神圣临在的经常光照，早已启示并向你显明的隐秘奇事；而是希望引导未受教育者进入光明，并在那些不知其所以然的人面前，显明你虔诚事工的原因和动机。\par

2. 诚然，你为至高天主每日在全世界的朝拜与尊崇所做的崇高努力，是普世赞美的主题。但是，你在我们巴勒斯坦行省、在那如同源泉般将救主圣言传于万民的城市中所奉献的感恩记录，以及你为纪念他战胜死亡而建造的圣殿与祝圣的教堂，还有那些崇高宏伟的建筑，为永远纪念救主坟墓而建的帝制丰碑——这些事的原因并非人人皆知。\par

3. 确实，那些因圣神的能力而得天上知识光照的人，明白其原因，并理当羡慕你、为你祝福，因那谋划与决心是上天所感。反之，无知和灵明黑暗的人则以公开的嘲笑与讥讽看待这些工程，认为这样一位伟大的君王对死者的坟墓和纪念碑如此热忱，实在奇怪而不足道。\par

4. \Quote{难道不是更好吗？}这样的人会说，\Quote{持守那些由古老传统尊崇的礼仪，寻求那些各省份都奉行的神明与英雄的恩宠，而不是因他们遭遇人性的苦难就拒绝和否认他们？他们当然可与那亲身受苦者同享尊荣；或者，若因他们不免于人类悲痛而应被拒绝，那么对他也可作同样的判决。}如此，他皱起重要而狭隘的眉头，以夸大的言辞，吐出他自以为是智慧的话。\par

5. 充满对无知者的怜悯，我们至慈之父的恩惠圣言，并非只邀请这样的人，而是邀请所有迷途者接受神圣知识的训导；并已在整个世界，在各乡各城，在文明与荒野之地，在每座城市，设立了这种训导的方式；作为一位恩惠的救主和灵魂的医师，他号召希腊人与外邦人、智者与无知者、富人与穷人、仆人与主人、臣民与君主、不敬者、亵渎者、无知者、作恶者、谤渎者，都来亲近，并赶紧接受他的天上医治。昔日他曾清楚地向众人宣布先前过犯的赦免，说：\BibleQuote{凡劳苦和负重担的，你们都到我这里来，我要使你们安息。}\BibleRef{玛 11:28}又说：\BibleQuote{我来不是召义人，而是召罪人悔改。}他随后说明了原因，说：\BibleQuote{健康的人不需要医生，有病的人才需要。}\BibleRef{玛 11:12}又说：\Quote{我不愿罪人死，而愿他悔改。}\par

6. 因此，只有那些自身受过神圣事物教导、并理解这些工程所源自的热忱动机的人，才能领会引导我们皇帝的超越人力的冲动，钦佩他对天主的虔诚，并相信他对我们救主复活纪念碑的关怀是一种从上而来的渴望，确实是由那位君王所感动；成为他忠信的仆人与良善的管家，乃是他最引以为傲的事。\par

7. 既然如此，我深信你定会赞同，最伟大的皇帝，我愿在此刻向众人宣告你虔诚工作的各种理由和动机。我愿作为你意图的阐释者，解释一个献身于爱天主的灵魂的深思。我打算教导所有人，凡是关心理解我们救主天主如何运用其权能的人，都该知道：那预存的万有掌管者为何最终从天上降下给我们；他为何取了我们的本性，甚至屈服于死亡权下；我将宣告随之而来的不朽生命和祂从死者中复活的原因。再次，为了那些仍需如此见证的人，我将提出令人信服的证据和论证。\par

8. 现在，让我开始我指定的工作。\par

那些将唯独应归于创造并统治世界的天主的崇拜，转给祂手的受造物；那些将太阳、月亮，或这个物质系统的其他部分，甚至元素本身——土、水、空气、火——与它们的造物主同等尊崇；那些把神的名号给予那些若非服从于创造世界的天主圣言、根本不会有存在甚至名称的事物：在我看，这样的人就像那些忽略赋予王宫壮丽的主手，却惊叹于其屋顶和墙壁、装饰它们的多样美丽色彩和绘画，以及镀金的天花板和雕刻，将属于其作品的工匠的技巧赞美归于它们。而他们应将惊叹的原因，不是归于这些可见之物，而是归于建筑师本人，并承认技巧的证据确实显明，但唯独那位使它们成为现在这样者才拥有那技巧。\par

9. 再者，我们可将那些人比作儿童：他们赞美七弦琴，却忽略发明或能使用它的人；或那些忘记勇敢战士，用胜利的花冠装饰他的矛和盾的人；或者最后，那些将一座伟大王城的广场、街道、公共建筑、神庙和体育场与它的创立者同等尊崇的人，忘记了他们的赞美应当归于不是无生命的石头，而是用智慧设计并执行这些伟大工程的那一位。\par

10. 那些用肉眼观察这个宇宙，并将其起源归诸太阳、月亮或任何其他天体的人，同样荒谬。更让他们承认这些本身就是更高智慧的工，记住它们的制造者和塑造者，并将赞美和尊崇归于祂，超越所有受造之物。更确切地说，受这些景象本身的激励，让他们以全心立志去荣耀和朝拜那一位：祂现在对肉眼不可见，但被灵魂清明无云的视觉所感知——那至高的、君临万有的天主圣言。以人的身体为例：没有人曾将智慧的属性归于智者或学者的眼睛、头、手、脚或其他成员，更不用说他的外衣了；没有人称哲学家的家具和器皿为智慧；但每个有理性的人都赞美那不可见而隐藏的力量——那人自身的心灵。\par

11. 那么，我们的赞美更应归于谁呢？不是归于宇宙的可见机制——它是物质的，由相同的元素构成；而是归于那不可见的圣言，祂塑造并安排了这一切，祂是天主的独生子，那超越万有的万有制造者从自身生出了祂，并立祂为这个宇宙的主宰和管理者。\par

12. 因为既然有死的身体或祂所创造的理性精神，因其与祂的完美相隔不可测度的距离，不可能接近至高天主——因为祂是未受生的，超越并凌驾于一切受造之上，不可言喻，不可接近，无法企及，正如祂的圣言所保证的，居住在\BibleRef{弟前 6:16}无人能近的光中；但它们是从无中受造的，与祂的未受生本质相距无限遥远；所以，全恩宠的全能天主在祂自己与它们之间设置了一种中间的力量，即祂独生圣言的神圣全能。这力量与父完美地亲近和结合，父居于祂内并分享祂的秘密计划，却仍因恩宠的满全，屈尊就卑，仿佛使自己适应那些远离天主至高威严者。否则，那超越万有者如何能按照祂自己的圣洁与可朽坏的有形物质结合呢？因此，圣言将自己与这个宇宙连接，将世界的缰绳握在手中，如同一位熟练的驭手，按照自己的旨意和喜悦转动和引导它。\par

13. 这些论断的证据是明显的。因为假设世界的这些组成部分（我们称之为元素，如地、水、气、火），其本性显然没有理智，是自存的；并且它们有一个共同的本质，精通自然科学的人称之为万物的伟大容器、母亲和乳母；而这个本质本身完全无形无相，没有灵魂和理性；那么，我们该说它从何获得现今的形式和美丽呢？我们将元素的区分或本性相反之物的结合归因于谁？谁命令液态水支撑住重元素土？谁阻止了水的下流，使它们升到云端？谁约束了火的力量，使它潜藏在木材中，并与其本性最相反的物质结合？谁将冷气与热混合，从而调和了对立原则的敌对？谁设计了人类种族的连续繁衍，并赋予它一种无尽的持续期限？谁塑造了男女形态，完美和谐地调整了他们的相互关系，并给予每个生物一个共同的生殖原理？谁改变了流动易朽、本身缺乏理性的种子的特性，并赋予它繁殖的力量？谁此刻正在运作这些和比这些更奇妙的万千效果——甚至超越一切奇妙——并以无形的影响日复一日地促使这一切产生？\par

14. 的确，行奇事的、真正全能的圣言可被视为这一切的动因：这圣言散布于一切受造中，以无体的能力渗透高处和深处，用其有力的掌握拥抱宇宙的长度和广度，整编并整顿了整个系统，从它无理性、无形的质料中为自己塑造了一件完美和谐的乐器，以全智无误的技巧弹奏其平衡的琴弦和音符。祂以不可解的规律治理太阳、月亮和天上其他的发光体，并为整体的服务指引它们的运动。\par

15. 正是这圣言屈尊降世，居住在我们的大地上，创造了多种多样的动物和植物界的美丽种类。正是这同一位圣言深入了深渊的隐密处，赋予了鱼类存在，并产生了那里存在的无数生命形式。是祂塑造了胎中的负担，并在自然的工坊中以生命的原则赋予其形状。藉着祂，流动的重湿气升上高处，然后通过净化的变化变甜，以适当的数量降在地上，并按规定的季节更丰沛地供应。\par

16. 像一位熟练的农夫，祂充分灌溉土地，恰当地调和干湿，用鲜艳的花朵、多彩的美丽面貌、宜人的芬芳、交替的果实种类和无数满足人口味的享受来点缀万物。但是，我为何敢于尝试一项无望的任务，去叙述圣言的伟大工作，描述一种超越凡人思想的能力呢？的确，有些人称祂为宇宙的本性，另一些人为世界灵魂，还有一些人为命运。另有些人竟宣布祂就是至高的天主自己，奇怪地混淆了最不同的事物：将那位至高无上的、未受生的能力、万有的主，带到这类地上，与可朽坏的有形身体结合，并给祂一个介于无理性的动物和有理性的人与不朽者之间的位置。\par

% structure: p0138 source=h2 target=section reason=relative-heading-rank; base=h2

\section{第十二章}

1. 另一方面，神圣的教义教导说，那至高之善的源头和万有的原因，是无法理解的，因此无法用言语、讲话或名称表达；不仅超越语言的能力，而且超越思想本身。不受地方或身体限制；既不在天上，也不在以太空间，也不在宇宙的任何其他部分；而是完全独立于一切其他事物，祂渗透了未探索的、隐秘的智慧的深处。神圣的启示教导我们承认祂为唯一真天主，脱离一切有形本质，区别于所有从属的职事。因此，据说万物来自祂，而非藉着祂。\par

2. 他本身以君王的身份居住在不可接近的光明的隐秘与未发现的领域，以他自己旨意的唯一权能命令并安排万有。凡存在的，皆因他的旨意而存在；离了他的旨意，便不能存在。他的旨意总是为了善，因为他本身就是善。然而万物藉着他而存在，即天主圣言，他以不可言喻的方式从至高的圣父——如同从永远不竭的泉源——发出，如同一条满溢着丰沛权能的河流，为保全整个宇宙而奔流。\par

3. 现在让我们从自身的经验中选取一个比喻。我们内在那无形且未被发现的理智，其本质无人曾知，它如同一位君主隐于内室，独自决断我们的行动方针。从这理智发出\Quote{独生子之言}出自其父之怀，以我们无法理解的方式和权能受生，它是其父思想的第一使者，宣告其隐秘的谋划，并传达到他人的耳中，成就其设计。\par

4. 因此，这种官能的益处为众人所享；然而，从未有人看见那无形且隐藏的理智，它是言语本身的父母。同理，更确切地说，以远超一切相似或比拟的方式，至高天主的完美圣言，作为圣父的独生子（并非在于发声的能力，也不被音节和词类所限，也不藉着在空气中振动的嗓音传达；而是他本身就是永生天主活而有效的圣言，并位格性地存有，是天主的德能与智慧），从他的圣父的天主性与国度发出。因此，作为完美圣父的完美后裔，万有的共同保存者，他以活生生的德能遍及受造界，从他的圆满中涌流出丰沛的理性、智慧、光明及一切其他福祉，不仅赐予离他最近的对象，也赐予最远的，无论是在地上、海上或任何其他领域。\par

5. 他以至公的公平为这一切设定界限、处所、法律与继承，按他君王的旨意分配给各人合适的部分。他赐予一些人超地的区域，另一些人则以天堂本身为居所；他将另一些人安置在以太空间，有些在空气中，有些仍在地上。正是他将人类从此处迁往另一境域，公正地审查他们在此的行为，并按各人的生平和习惯给予赏报。藉着他，不仅理性的受造物，连为服务人而造的牲畜，其生命与食物也都得到供给；\par

6. 对后者，他赐予可朽坏且短暂存有的享受；对前者，他邀请他们分享不朽生命的拥有。天主圣言的作用何其普遍：他无处不在，以他的智力的权能渗透万有，仰望他的圣父，依照他的旨意治理这低级受造界——这是次于他并源于他的，作为万有的共同保存者。\par

7. 这圣言仿佛是居间的，吸引受造者趋向未受造的本质，作为两者之间不可断裂的纽带，以不可分离的绳索联结最遥远的事物。他就是治理宇宙的天主眷顾；全世界的护卫与指导：他是天主这位独生天主的德能与智慧，是天主自己所生的圣言。因为\BibleQuote{在起初已有圣言，圣言与天主同在，圣言就是天主。万有是藉着他而造成的，凡受造的，没有一样不是藉着他造成的}；正如我们从圣史的话语所学到的。\BibleRef{若 1:1} 藉着他赋予生命的力量，整个自然界生长繁茂，被他持续的甘霖所滋润，并被赋予永不枯竭的活力与美丽。\par

8. 他引导宇宙的缰绳，使其运行符合圣父的旨意，并好像转动这艘巨轮的船舵。这荣耀的代理者，至高天主的独生子，由圣父所生，作为他的完美后裔，圣父把他赐给这世界作为最高的恩惠；将他的圣言如同精神注入无生命的身体，注入无意识的自然；藉着神圣的权能，将光明与活力赋予那原本粗糙、无生命、无定形的物质。因此，我们应当承认并敬仰他无处不在，赋予物质与自然元素生命：在他内我们看见光明，即那不可言喻之光的属灵后裔：在本质上是一，因为他是唯一圣父之子；但他自身拥有许多不同的权能。\par

9. 世界确实分成许多部分，但我们不可因此以为有许多独立的代理者；即使受造工程繁多，我们也不可由此假设有许多神祇。那些幼稚而痴迷的多神崇拜者的错误何其严重！他们神化宇宙的组成部分，将那唯一体系分割成许多！\par

10. 这种行为如同有人将一个人的眼睛抽象出来，称之为这个人自身，又把耳朵称为另一个人，头也如此；或凭思想的努力分离颈、胸与肩、脚与手或其他肢体，甚至感觉的能力，从而宣称一个人是一群人。这种愚行必受有识之士的鄙弃。然而，从单一世界的组成部分为自己臆想出众多神祇，甚至认为那由创造者所造、由许多部分构成的世界本身就是神，正是如此：他不知道神圣本质绝对不能分割成部分；因为如果是复合的，就必须藉另一个权能的作用；而复合的东西绝不能是神圣的。它怎么可能如此，若由不等同且不相似的、因而有优劣的元素组成？神圣本质是单纯、不可分、非复合的，存在于远远超越这世界可见构造的无限高处。\par

11. 因此，我们由神圣的宣告者的明确见证确知，那先在万有之先的天主圣言，必须是一切有灵之物的唯一保存者；而那超越万有、为圣言产生之根源的天主，祂自身既是万有的原因，理当被称为圣言之父，如对祂的独生子，祂自己不承认有任何更高的原因。所以，天主自身是独一的，从祂发出那唯一的独生圣言，无所不在的万有保存者。正如多弦的里拉琴由不同音调的和弦组成，有高有低，有些松驰，有些紧张，有些介于两个极端之间，然而全都按照和声艺术的规则调谐；同样，这个物质世界，由许多元素组成，包含相反和对抗的原理，如润湿与干燥、寒冷与炎热，却浑然成为和谐的整体，正可称为天主手所造的伟大乐器：一件乐器，由那圣言——祂自身不是由部分或对立原理组成，而是不可分割、单纯无杂——以完美的技巧演奏，产生出一首既符合祂的圣父、万有至高主宰的旨意，又对自己荣耀的旋律。再者，正如在一个身体内包含许多外在和内在的肢体与部分，却有一个不可见的灵魂、一个不可分割和非物质的理智贯穿全体；同样，在这个受造界中，它由许多部分组成，却只是一个：同样，那伟大而全能的天主圣言，贯穿万有，以坚定不移的能力遍布宇宙，是其中一切存在事物的原因。\par

12. 你环顾这有形世界的范围。难道不见同一穹苍包含星辰无数的轨道和群吗？再者，太阳只有一个，却以其光芒的卓越光辉遮掩了许多，不，是遮掩了所有其它发光体。同样，正如圣父是独一的，祂的圣言也是独一的，那完美圣父的完美圣子。若有人因它们不是更多而反对，同样可以抱怨没有多个太阳、月亮或世界，以及上千种其它事物；如同那疯狂的人，想要颠覆大自然本身美好而完美的进程。正如在有形的世界，同样也在精神的境界：在前者，同一太阳将其光遍布这个物质大地；在后者，那唯一全能的天主圣言以无形隐秘的大能光照一切。\par

13. 再者，在人内有一个灵和一个理智官能，它却是无数效果的动力因。同一个心智，受教于许多事物，会尝试耕种土地、建造并驾驶船只、建造房屋：是的，人的同一心智和理性能够以千百种形式获取知识：同一个心智会理解几何学和天文学，并讨论文法、修辞和医术的规则。它不仅在科学上卓越，也在实践上如此：然而，无人曾假设在一个人类形体中存在许多心智，也未曾因人类具有如此多样的知识能力而惊叹于其存在的多元性。\par

14. 假如有人找到一块无定形的泥土，用手捏塑它，赋予它一个生物的形像；头形一个样，手足另一个样，眼睛和脸颊又一个样，并照此造出耳朵、口鼻、胸肩，遵照造型艺术的规则。结果，在一个身体上确实有形状、部分和肢体的多样性；然而我们不能认为这是众多手的作品，而应完全归功于一位艺术家的技巧，并将赞颂归于他，因他凭单一心智的能力塑造了这一切。宇宙本身也是如此，它虽由许多部分组成，却是一体的：但我们确实无需假设许多创造能力，也无需发明多神。我们的职责是钦崇那全智全美的行动者，祂实在是天主的德能和智慧，祂那不可分割的力量和能力弥漫并渗透宇宙，创造并赋予万有生命，并为整体和个体提供那丰富的供给，而祂自己就是这供给的源泉。\par

15. 同样，太阳光线同一印象同时照亮空气，给予眼睛光明，给触觉温暖，给大地丰饶，给植物生长。同一发光体构成时间的过程，治理星辰的运行，执行天穹的周行，赋予大地美丽，并向众人显示天主的大能：所有这一切，它都凭其本性的独一而不假外求的能力完成。同样，火有精炼金子、熔化铅、溶解蜡、烘干泥土、烧毁木头的性质；凭同一燃烧的能力产生这些不同的效果。\par

16. 同样，至高的天主圣言，贯穿万有，无处不在，在天地间处处临在，治理并引导有形无形的受造界——太阳、穹苍和宇宙本身——祂的能力在本性上不可言喻，在效果上不可抗拒。从祂，如同从永泉，太阳、月亮和星辰接受光明：祂永远统治那祂所塑造作为祂伟大恰当象征的穹苍。超越天地领域的天使和神体、无形无体的有灵之物，被祂以祂独有的宝藏充满光明和生命、智慧和德能、一切伟大和美好。再者，以同一创造力，祂不断赋予元素实体，调节结合与组合、形状与形象、以及有机体的无数性质；保持动物和植物生命、理性与无理性受造物的不同区别；并以同等的能力供给万有一切：从而证明祂自己是创造者，不是里拉琴，而是那完美和谐的系统，那是由那独一创造世界的圣言的作品。\par

% structure: p0155 source=h2 target=section reason=relative-heading-rank; base=h2

\section{第十三章}

1. 现在我们继续解释天主的圣言\OriginalTerm{天主的圣言}{Word of God}为何如此强大，竟降来与人同居。我们这无知愚昧的种族，不能理解那天地之主，祂从圣父的性体流出，如同从至高源头涌出，遍临世界各处，并以最清晰的证据显明祂对人类福祉的上智安排；我们却将那当受钦崇的天主尊号，归于太阳、月亮、天穹和天上的星辰。不仅如此，还神化了大地本身、大地的出产以及维持动物生命所需的各种物质，并为刻瑞斯\OriginalTerm{刻瑞斯}{Ceres}、普洛塞庇娜\OriginalTerm{普洛塞庇娜}{Proserpine}、巴克斯\OriginalTerm{巴克斯}{Bacchus}以及许多此类神祇塑造了形象。\par

2. 不仅如此，他们甚至不羞于将自己内心的意念和表达意念的言语称为神：称心灵本身为密涅瓦\OriginalTerm{密涅瓦}{Minerva}，言语为墨丘利\OriginalTerm{墨丘利}{Mercury}，并将谟涅摩绪涅\OriginalTerm{谟涅摩绪涅}{Mnemosyne}和缪斯\OriginalTerm{缪斯}{Muses}之名赋予那些借以获得学识的能力。这还不够；他们更在亵渎与愚妄的道路上愈行愈远，神化了自己邪恶的情欲——这本应令他们憎恶，或用自制之德加以约束。他们将自己淫欲、情欲、灵魂的不洁之疾、引人淫秽的身体肢体，乃至在可耻享乐中毫无节制的情状，都冠以丘比特\OriginalTerm{丘比特}{Cupid}、普里阿普斯\OriginalTerm{普里阿普斯}{Priapus}、维纳斯\OriginalTerm{维纳斯}{Venus}以及其他类似的名称。\par

3. 他们甚至没有止步于此。他们将关于天主的意念降低到这个有肉身和必死的生命，把自己的同胞奉为神明，将神和英雄的名号加诸那些经历了众人共同命运的人身上，虚妄地以为那神圣不朽的本体能常临于死者的坟墓和纪念碑之间。不仅如此：他们还向各种飞禽走兽乃至最毒恶的爬虫献上敬神之礼；砍伐树木，开凿岩石；为自己预备铜、铁及其他金属，用它们塑造男女形象、走兽和爬虫的形像，并将这些作为崇拜的对象。\par

4. 这还不够。他们甚至将同样的敬神之礼归于那些潜伏在雕像之内、或藏于隐秘幽暗之处、渴望吸饮祭酒和嗅取祭品烟气的邪灵自身。不仅如此，他们又借助禁咒的魔法和不洁的歌曲与咒语，试图获取这些邪灵以及在天空中流动的不可见力量的日常帮助。此外，不同民族采用不同人物作为敬拜对象。希腊人将巴克斯\OriginalTerm{巴克斯}{Bacchus}、赫拉克勒斯\OriginalTerm{赫拉克勒斯}{Hercules}、埃斯库拉庇乌斯\OriginalTerm{埃斯库拉庇乌斯}{Æsculapius}、阿波罗\OriginalTerm{阿波罗}{Apollo}及其他本是凡人者冠以神和英雄的称号。埃及人则将荷鲁斯\OriginalTerm{荷鲁斯}{Horus}、伊西斯\OriginalTerm{伊西斯}{Isis}、奥西里斯\OriginalTerm{奥西里斯}{Osiris}及其他类似凡人奉为神明。于是，那些夸耀自己凭奇妙技艺发现了几何学、天文学和数学的人，虽在自己的智慧中自以为是，却不知道也不懂得如何估量天主之能的尺度，或计算祂那远超越无理性的必死之物的至高伟大。\par

5. 因此，他们不羞于将神的名号加在最丑陋的受造物、毒虫和猛兽之上。腓尼基人尊崇麦勒卡特\OriginalTerm{麦勒卡特}{Melcatharus}、乌索鲁斯\OriginalTerm{乌索鲁斯}{Usorus}及其他不过是凡人、且几乎不值得尊重的人；阿拉伯人尊崇杜萨里斯\OriginalTerm{杜萨里斯}{Dusaris}和奥博达斯\OriginalTerm{奥博达斯}{Obodas}；盖塔人尊崇札摩克西斯\OriginalTerm{札摩克西斯}{Zamolxis}；奇里乞亚人尊崇摩普索斯\OriginalTerm{摩普索斯}{Mopsus}；底比斯人尊崇安菲阿拉俄斯\OriginalTerm{安菲阿拉俄斯}{Amphiaraus}；简言之，每个民族都有自己独特的神祇，这些神祇与他们的凡人同胞毫无区别，纯粹是真实的人。再者，埃及人一致同意，还有腓尼基人、希腊人，乃至天下万民，全都一同敬拜世界的组成部分和元素，甚至地上出产的东西。最令人惊讶的是，尽管他们承认这些神祇犯有奸淫、违反伦常和放荡的罪行，他们不仅用神庙、神龛和雕像充满每座城市、村庄和地区，还效法这些邪恶榜样，自毁灵魂。\par

6. 我们听说他们所谓的神和神子被描述为英雄和善灵，这些称号与真理全然相悖，这些尊荣与它们所要夸耀的品性完全不符。就好像有人想要指出太阳和天上的明灯，却不愿举目仰望，反而用手在地上摸索，在污泥和烂泥中寻找天上的力量。同样，人类被自己的愚妄和邪灵的诡计所欺骗，竟相信那远超天地之上的神圣属灵的本体，能与必死的身体在此世有出生、情感和死亡。他们疯狂到如此地步，竟将最亲爱的亲人祭献给他们所谓的神，不顾一切天然联系，在疯狂冲动的驱使下杀死他们唯一的、最心爱的子女。\par

7. 因为还有什么比将人祭献、以同族人的血玷污每座城市乃至自己的房屋更显疯狂的呢？希腊人自己不也证实了这一点，且全部历史不都充满了同样不虔的记载吗？腓尼基人将他们最心爱的独生子女作为年祭献给\OriginalTerm{萨图尔努斯}{Saturn}。罗得岛人在梅塔格特尼翁月的第六日，将人牲献给同一神祇。在萨拉米斯，一个人在\OriginalTerm{弥涅耳瓦·阿格劳莉丝}{Minerva Agraulis}和狄俄墨得斯的神庙中被追逐，被迫绕祭台跑三圈，随后被司铎用长矛刺穿，放在熊熊火堆上焚烧作为全燔祭。在埃及，人祭最为盛行。在赫利奥波利斯，每日有三名牺牲献给\OriginalTerm{朱诺}{Juno}，阿莫西斯王对这种残忍行为深感震惊，下令代之以同等数量的蜡像。在希俄斯，以及后来在忒涅多斯，一个人被杀死并献给\OriginalTerm{奥马迪·巴克斯}{Omadian Bacchus}。在斯巴达，他们将人祭献给\OriginalTerm{玛尔斯}{Mars}。在克里特，他们同样将人祭献给萨图尔努斯。在叙利亚的劳迪西亚，每年有一名童贞女被杀死以尊崇\OriginalTerm{弥涅耳瓦}{Minerva}，如今则以一头牡鹿代替。利比亚人和迦太基人以人牲平息神怒。阿拉伯的杜马泰尼人每年将一个男孩埋在祭台下。历史告诉我们，希腊人无例外，还有色雷斯人和斯基泰人，在出征前都惯常进行人祭。雅典人记载了劳斯的童贞女儿和厄瑞克透斯的女儿被献祭。谁不知道，直到今天，在罗马的\OriginalTerm{朱庇特·拉提亚里斯}{Jupiter Latiaris}节上仍有人牲被献？\par

8. 这些事实得到了最权威哲学家的证明。图书编纂者狄奥多罗斯断言，利比亚人曾将两百名最高贵的青年祭献给萨图尔努斯，另有三百名由他们自己的父母自愿献出。罗马史编纂者狄奥尼修斯明确说，\OriginalTerm{朱庇特}{Jupiter}和\OriginalTerm{阿波罗}{Apollo}在意大利向所谓的原住民要求人祭。他记述说，为此他们献出各种出产的一部分给神祇；但因拒绝杀人祭献，他们陷入了多种灾祸，直到他们杀死了十分之一的人口才得解脱，这次生命的祭献导致了其国家的荒凉。如此巨大的祸患曾经折磨着全人类。\par

9. 这还不是他们痛苦的极限：他们还呻吟在同样众多且无可补救的其他祸患的重压下。普天下所有民族，无论文明或野蛮，都仿佛被恶魔般的狂乱所驱使，感染上叛乱，如同某种凶猛可怖的疾病；以致人类家族彼此无可挽回地分裂；社会的大系统支离破碎；在地球的每一个角落，人们彼此对立，为法律和政府的问题激烈争斗。\par

10. 更有甚者：他们被激起的愤怒驱使，彼此陷入冲突，如此频繁以致他们的生命就像在不间断的战争中度过。没有人能在不准备迎敌的情况下出行；在乡间和村庄，乡下人佩带刀剑，配备甲胄而非农具，并以掠夺和奴役任何邻国人为高尚功业。\par

11. 更有甚者：从他们自己为神祇编造的寓言中，他们推导出卑鄙放纵生活的借口，并以各种放荡行为毁坏身体和灵魂。不仅如此，他们甚至越过了自然所定的界限，共同犯下了不可想象且不可名状的罪行，\Quote{男人与男人（如圣作者所言）行可耻之事，并在自己身上受到他们错谬所应得的报应。}\par

12. 他们甚至没有止步于此；反倒歪曲了对天主的天然思想，否认世界的进程由祂的眷顾所引导，而将万物的存在与构造归于盲目的机缘或命运的必然。\par

13. 再者：他们相信灵魂与身体皆因死亡而消散，于是过着兽性的、不配其名的生活；对灵魂的本质或存在漠不关心，不畏惧神圣公义的审判台，不期望德行的赏报，也不将惩罚视为恶行的报应。\par

14. 因此，整个民族成为各种形式邪恶的掠食对象，因自身的野蛮行为而衰败：有些人与母亲、与姐妹行最卑劣无耻的乱伦，还有人玷污自己的女儿。有人发现他们杀害信任的宾客；有人吃人肉；有人勒死老年人并宴食之；有人将他们活活扔给狗。若要试图描述那已统治全人类的根深蒂固之恶疾的百般症状，我的时间是不够的。\par

15. 这些，以及无数其他类似的，是普遍存在的邪恶，为此，恩慈的天主圣言、对祂的人类羊群充满怜悯，早就通过祂先知的服务，以及更早和更晚通过那些以虔诚献身于天主而著称之人的服务，邀请那些绝望受苦者来医治自己；并藉着法律、劝诫和各种教义，向人类宣告真正虔敬的原则和要素。但是，当人类如我所说被撕裂和搅乱，不是被豺狼和野兽，而是被无情且毁灭灵魂的恶灵所害，人的能力已不足，需要超乎人的帮助时，天主圣言服从祂全恩慈圣父的旨意，终于亲自显现，并极其情愿地居于我们中间。\par

16. 我已讲述了祂来临的原因，促使祂屈尊与人类交往；祂并非以惯常的形式和方式，因为祂是无形的，且遍及全世界，以祂在天地间的作为证明祂全能之力的伟大，而是以一种新的、前所未闻的方式。祂取了一个可朽的身体，屈尊与人交往交谈；愿意借着与他们相似的形体，拯救我们有死的人类。\par

% structure: p0172 source=h2 target=section reason=relative-heading-rank; base=h2

\section{第十四章}

1. 现在让我们解释：为何无形体的\OriginalTerm{圣言}{Word of God}取了这个可朽的身体作为与人类交往的中介。因为若非以人形，那神圣而不可触摸、非物质而\OriginalTerm{不可见的性体}{invisible Essence}，如何能向那些在受造物和尘世事物中寻求天主、却无法或不愿以其他方式认识万物的创造者和制造者的人显现自己呢？\par

2. 因此，作为与人类交往的适宜方式，祂取了一个可朽的身体，因为这是人类所熟悉的；正如谚语所说，\Quote{相似者爱相似者}；对于那些眷恋可见之物、在雕像和无生命的偶像中寻找神祇、想象神性存在于物质和形体的实体中、甚至将神性之名赋予世人的人，\OriginalTerm{圣言}{Word of God}便以这种形象出现在他们面前。\par

3. 因此，祂为自己预备了这个身体，作为\OriginalTerm{三重圣化的殿宇}{thrice-hallowed temple}，一个智能之力的可感居所；一个高贵而最神圣的\OriginalTerm{形像}{form}，远胜于任何无生命的雕像。那由低微工匠之手用铜、铁、金、象牙、木或石塑造的物质而无感觉的偶像，或许是邪灵的适宜居所；但那由天上智慧之力所造的\OriginalTerm{神圣形像}{Divine form}，则拥有生命和属神的存在；这是一个被各种美德所充满的形像，圣言的居所，圣天主的圣殿。\par

4. 如此，内住的圣言与人类交谈，被人认识，视为与他们相同；然而祂并不屈服于他们那样的情欲，也不像自然灵魂那样受身体支配。祂未失去任何本有的伟大，也未改变祂的\OriginalTerm{本有神性}{proper Deity}。因为正如太阳的\OriginalTerm{遍照的光辉}{all-pervading radiance}接触死亡和不洁的身体而不受玷污；同样，圣言的无形体之力也绝不会因与人类身体的神性接触而损害其\OriginalTerm{本质的纯洁}{essential purity}，或失去任何伟大。\par

5. 因此，我所说的，我们\OriginalTerm{共同的救主}{common Saviour}证明自己是所有人的恩主和保护者，借着祂人性的工具彰显祂的智慧，如同乐师用竖琴展示其技巧。希腊神话告诉我们，\Person{俄耳甫斯}能以大师之手拨动琴弦，驯服凶猛的野兽，平息其野性；希腊人传颂并普遍相信这一故事：无意识的乐器能制服未驯服的野兽，并因其悦耳之力而让树木移动。但那位完美和谐的创造者、\OriginalTerm{全智的圣言}{all-wise Word of God}，渴望为人类灵魂的\OriginalTerm{多种疾病}{manifold diseases}施治各种药方，便使用祂智慧所造的人性作为乐器，以它的悦耳之声安抚——并非野兽，而是有理智的蛮族；借着祂神圣教义的治愈之力，治愈了文明和野蛮民族中每一种暴躁的性情、每一种狂怒而愤怒的灵魂之情欲。如同一位技艺精湛的医师，祂以合适而类似的药方，应对那些在自然和形体中寻求天主的灵魂的疾病，并以人形向他们显示天主。\par

6. 然后，祂对身体和灵魂同样关怀，在人眼前呈现\OriginalTerm{奇事与征兆}{wonders and signs}，作为祂神性能力的证明，同时将祂用有形的舌头所宣讲的教义注入他们肉体的耳中。简而言之，祂借着为那些无法认识祂神性的人所取的身体，完成了祂的一切工程。\par

7. 在这一切中，祂是祂父旨意的仆人，祂本身仍与在父那里时相同：\OriginalTerm{本质不变}{unchanged in essence}，\OriginalTerm{本性无损}{unimpaired in nature}，不被可朽血肉的束缚所困，也不因居住在人类身体中而无法临在他处。\par

8. 不仅如此，在祂与人交往的同时，祂\OriginalTerm{遍及万物}{pervading all things}，与父同在并内在于父，甚至在那时也照顾天地间的一切。祂不像我们一样被限制在某个地方，或被阻止持续运用祂的神性能力。祂将自己的给予人，却不收取任何回报；祂将神性能力赋予必朽者，却未从必朽性本身获得任何增益。\par

9. 因此，祂的人性诞生并未给祂带来任何玷污；祂那\OriginalTerm{不动情的性体}{impassible Essence}也未曾因祂可朽身体的解体而受苦。因为假设一架竖琴意外受损，或琴弦断裂，演奏者并不会因此受苦；同样，若智者的身体遭受刑罚，我们不能断言他的智慧或内在的灵魂受损或被焚烧。\par

10. 我们更不可能断言\OriginalTerm{圣言内有的能力}{inherent power of the Word}因祂\OriginalTerm{身体的苦难}{bodily passion}而遭受任何损害，正如我们已使用的例子：从天上射向地面的\OriginalTerm{太阳光线}{solar rays}接触污秽和各种污染物，却不玷污。我们确实可以说这些事物分有光的照耀，但光并未被污染，太阳也未因与其他物体的接触而受玷污。\par

11. 再者，这些事物本身并不违反本性；但救主、\OriginalTerm{无形体的圣言}{incorporeal Word of God}既然是\OriginalTerm{生命与属神之光}{Life and spiritual Light}本身，凡祂以神圣而无形体之力所触摸的，必然获得光与生命的智能。因此，若祂触摸一个身体，它便被光照和圣化，立即脱离一切疾病、软弱和痛苦，之前所缺乏的由祂的\OriginalTerm{圆满}{fullness}的一部分所补充。\par

12. 他在地上的生活就是如此；时而在人的本性上证明他与我们同感，时而显现为天主的圣言：作为天主，他的作为奇妙而大能；预言遥远的未来；在每一行动中，以征兆、奇迹和超然能力宣告那少为人知的圣言；最终，通过他的神圣教导，邀请人的灵魂预备那超乎诸天的住所。\par

% structure: p0185 source=h2 target=section reason=relative-heading-rank; base=h2

\section{第十五章}

1. 如今还剩下什么，不过是陈述那一切的高峰事实：即他的死亡（如此广为人知）、他受难的方式，以及死后复活的大奇迹；然后以最清晰的见证确立这些事件的真相。\par

2. 由于上述原因，他使用了可朽之身作为工具，作为与其神圣尊威相称的形像，如同一位大能的君王，用它作为与人交往的通译，一切作为都与他自身的神能相一致。那么，假设在他寄居人间之末，他若以任何其他方式突然从人眼前消失，并秘密除掉那为他通译的（即他所取的形像），急忙逃避死亡，然后自己使他的可朽之身归于腐朽和分解：无疑，在这种情况下，他会被众人视为纯粹的幻影。而且，若他——那生命、圣言和天主的能力——竟将这为他通译的交付给腐朽和死亡，那也不符合他的身份。\par

3. 同样，他与恶神的争战也不会通过与死亡之能的冲突而达到圆满。他退隐之处必然不为人知；未曾亲眼见他的人也不会相信他的存在。没有证据证明他胜过了死亡，他也不会使可朽者脱离其自然软弱的律法。他的名字从未传遍世界；他也不能使他的门徒蔑视死亡，或鼓励接受他道理的人盼望将来在天主内享受生命。他也不会实现自己应许的保证，也不会成就先知关于他的预言。他也不会经历那最后的冲突；因为这正是与死亡之能的搏斗。\par

4. 因此，基于所有这些理由，既然那为圣言服务过的可朽之身必须有一个与其神圣居住者相称的结局，死亡的方式便相应地被规定了。因为只剩下两种选择：要么将他的身体完全交付腐朽，从而使生命场景以不光彩的方式结束；要么证明他自己战胜了死亡，并以神圣大能的行为使可朽者成为不朽；前者会违背他自己的应许。因为正如火的本性不是冷却，光的本性不是黑暗，同样，生命剥夺生命是不相称的，神圣理智做不合理的事也是不相称的。因为，既然他应许给予别人生命，却容许自己的身体（他所拣选的形像）在腐朽之权下灭亡，这如何合理？他激发门徒对不朽的希望，却使这神圣旨意的解释者被死亡毁灭，这如何合理？\par

5. 因此，第二种选择是必要的：即他应当宣告自己对死亡之能的统治。但如何宣告？是隐秘而秘密地，还是公开地在众人眼前？如此伟大的成就，若是不为人知、未曾启示，则必对人无益；而同一件事，若被公开宣告并为众人理解，则会因其奇妙性质而惠及所有人。所以，既然有必要证明他的身体战胜了死亡，且不是秘密地而是在人眼前，他便没有退缩，因为那确实会显出恐惧和对死亡之能的卑下，而是维持了那使可朽者成为不朽的与仇敌的争战；这争战是为了所有人的生命、不朽和救恩。\par

6. 假设有人想向我们证明一个容器能抵御火力，他如何能更好地证明这事实？岂不是将它投入熔炉，然后完好无损地取出吗？同样，作为万有生命之源的圣言，想要证明那为人的救恩而取的身体战胜了死亡，并使这身体分享他自己的生命和不朽，便采取了与此目标一致的做法。他将自己的身体暂时留下，交付死亡以证明其可朽性，然后又迅速从死亡中赎回它，以证实那神圣的大能，藉此他显明了所应许的不朽完全超乎死亡的范围。\par

7. 原因很清楚。有必要让他的门徒亲眼证实那复活的可靠性，他曾教导他们将希望寄托于复活，以此作为超越死亡恐惧的动力。这真理的清晰印象对于立志度虔敬生活的人而言是非常必要的；对于那注定要普天下传扬他名、并传达天主知识（这知识他早已为万民预定）的人而言，则更为必要。\par

8. 因为这样的人需要对未来生命有最坚定的信念，以便能以无畏和不退缩的热忱，持守与外邦人和多神论错误的争战；这争战的危险，除非他们习惯于蔑视死亡，否则永远无法准备面对。因此，在武装门徒对抗这最后仇敌之权柄时，他不仅以言语规条传授教义，也不以动听的或可能的论证试图证明灵魂的不朽，而是亲自向他们展示了对死亡的真正胜利。\par

9. 这便是我们的救主与死亡权势争战的首要且最大的原因；借此他向门徒证明，那使全人类恐惧之物原是虚无，并提供了他所应许之生命的切实可见证据；仿佛呈献了我们共同希望——即在天主面前的未来生命与不朽——的初果。\par

10. 他复活的第二个原因，是要彰显那居于他可朽身体内的天主性大能。在此之前，人类曾向屈服于死亡权势的人奉献神圣尊荣，将神与英雄的称号授予与自己一样的凡人。因此，天主的圣言彰显了他的恩慈品性，向人证明自己超越死亡，召唤他的可朽身体归回第二次生命，在众人眼前展示对死亡的不朽胜利，并教导他们承认这一胜利的创始者即使在死亡中也是唯一真天主。\par

11. 我还可以提出救主死亡的第三个原因。他是为全人类而献于宇宙至高君王的牺牲：一个为人类所需而祝圣、并为推翻敬拜魔鬼之谬误而献上的牺牲。因为一旦那唯一神圣而强大的祭献——我们救主的圣身——为人而被宰杀，作为万民的赎价（他们此前陷于不虔迷信的罪债中），从此不洁与不圣的邪灵权势被彻底废除，一切生于尘世且引人迷惑的谬误立即被削弱并消灭。\par

12. 如此，这出自他们自身、有益于人的牺牲——我指的是圣言的可朽身体——是为全人类的共同族类而献上的。这就是交付于死亡的祭献，神圣圣经论及它说：\BibleQuote{请看天主的羔羊，除免世罪者。}\BibleRef{若 1:29} 又说：\BibleQuote{祂被牵去如羊待宰，又如羔羊在剪毛者前缄默。} 他们也说明了原因，说：\BibleQuote{祂承担了我们的罪，为我们而受苦：我们却以为祂在受罚、受苦、受难。但祂是因我们的过犯而被刺透，因我们的罪恶而被粉碎：使我们得平安的惩戒落在祂身上，\Emph{而且}因祂的伤痕我们得治愈。我们众人都如羊走迷了路，各人偏行己路；主却把祂献为我们众人的罪。}\par

13. 这就是促使天主的圣言献上他人性身体的原因。但既然祂是伟大的大司祭，祝圣于至高上主与君王，因而超越牺牲——祂是天主的圣言、大能与智慧；祂很快从死亡掌握中召回自己的身体，将它呈献给圣父，作为我们共同救恩的初果，并树立这战利品，作为祂战胜死亡与撒殚以及废除人祭的明证，为全人类带来祝福。\par

% structure: p0199 source=h2 target=section reason=relative-heading-rank; base=h2

\section{第十六章}

1. 现在，我们该向前继续证明这些事；如果这些真理确需证明，如果证言有助于确证明显事实的确定性。然而，此处将提供这样的证言；愿人以留心与恩慈的耳朵领受。\par

2. 古时，地上的万民，即全人类，被分散以各种形式隶属于省、国、地方的政府，受制于种种王国与公侯。这种多样性的后果是战争与争斗、人口减少与被掳，在乡村与城市中无休止地肆虐。由此，产生了历史上无数事件：奸淫、抢夺妇女；由此产生了特洛伊的苦难，与所有民族皆知的古代悲剧。\par

3. 这些事的起源可公正地归咎于多神论谬误的迷惑。但当那赎回我们的工具——基督三度圣洁的身体（它证明自己超越一切撒殚的诡计，在言行上远离邪恶）——被高举起来，既为废除古代邪恶，又为象征祂对黑暗权势的胜利时，这些邪灵的势力立即被摧毁。多元的政体、暴政与共和、由此引发的围城与国土荒废，都不再存在；一位天主被宣告给全人类。\par

4. 同时，一个普世权力——罗马帝国——兴起并繁荣，而民族之间持久且不可和解的仇恨被消除了；正如对一位天主、一种宗教与救恩之道（即基督的教导）的认知被宣告给全人类，在同一时期，罗马帝国的整个统辖权归于一位君主，深沉的和平统治了全世界。因此，藉着同一天主的明确安排，两个祝福的根源——罗马帝国与基督虔敬的教导——一同兴起，为人类造福。\par

5. 因为在此之前，世界上的各国，如叙利亚、亚细亚、马其顿、埃及和阿拉伯，分别隶属于不同的统治者。犹太民族又曾在巴勒斯坦地建立其统治。这些民族，在各乡村、城市与地区，被一种疯狂的精神所驱动，陷入不断且流血的战争与冲突。但两个强大的力量——自此由一位君主统治的罗马帝国，以及基督宗教——从同一起点出发，降服并调和了这些争斗的要素。\par

6. 我们救主的大能立即摧毁了黑暗权势的众多政府与众多神祇，并向所有人——无论未开化或文明者——直推到地极，宣告天主自身的独一主权。同时，罗马帝国，由于众多政府的原因已被移除，便轻易征服了那些仍存的；其目标是将所有民族联合成一个和谐的整体；这一目标在很大程度上已实现，并且注定要更加完美地达成，甚至直到居人之地的终极征服，凭借这有益于救恩的教导，并藉着那助其顺利进展的天主性大能的援助。\par

7. 这确实是一件奇事，对于那些怀着对真理的热爱来审视这个问题，并且不愿挑剔这些恩赐的人而言。恶魔迷信的虚妄被证实了；各民族之间根深蒂固的纷争和相互仇恨被消除了；同时，独一的天主和对天主的认识被宣告给所有人；一个普世的帝国盛行；全人类被和平与和谐的掌控力量所征服，彼此视如兄弟，并回应了他们共同本性的情感。因此，作为独一天主和圣父的子女，并承认真宗教为共同的母亲，他们以和平的话语互相问候和欢迎。如此，整个世界就像一个井井有条、团结一致的大家庭；每个人都可以无阻碍地旅行到他想去的任何地方；人们可以从西到东、从东到西安全地旅行，如同回到自己的故乡；简言之，古代先知的预言和预言的应验，其数量之多我们无法在此一一列举，尤其是那些关于救恩圣言所说的预言：\BibleQuote{祂必统治从这海到那海，从大河直到地极。} 再者，\BibleQuote{在祂的岁月中，正义要兴起，满盈和平。} \BibleQuote{他们将要把刀剑铸成锄头，将枪矛制成镰刀；民族不再举刀攻击民族，他们也不再学习战争。}\par

8. 这些话，早在多个世纪前以希伯来语预言，在我们这个时代得到了有形的应验，由此清晰地确认了古代先知的见证。现在，如果你还渴望更充分的证据，那么不是通过言语，而是从事实本身来接受它。睁开你理智的眼睛；拓宽思想的门扉；暂停片刻，思考；如同另一个自己一样审问自己，并如此勤勉地审视事物的本质。世上任何时代有什么君王或王子，什么哲学家、立法者或先知，在文明或野蛮之地，达到了如此伟大的卓越高度——我不是说死后，而是仍然活着并充满大能之时——使得全人类都满口称赞他的名字？确实，除了我们唯一的救主，没有谁做到了这一点。当祂战胜死亡后，祂对门徒说了这话，并藉由事件实现了它，对他们说：\BibleQuote{你们要去，因我的名使万民成为门徒。} 正是祂明确保证，祂的福音必须传遍全世界，作为对万民的见证，并立刻证实了祂的话：因为不久之后，世界本身就被祂的教义所充满。\par

9. 那么，那些在我演讲开头挑剔的人将如何回应这点呢？因为目睹的见证力量确实胜过任何言语的论证。除了祂，还有谁以无形却强有力的手，像驱赶野兽一样，从人类社会中驱逐了那一直有害和破坏性的恶灵族群，这些恶灵曾经使所有民族成为他们的猎物，并通过他们的偶像的运作在人类中施行了许多欺骗？除了我们的救主，还有谁藉着呼求祂的名，并藉着通过祂向至高天主献上的真诚祈祷，赐给那些真诚、顺服地遵循祂自己所指定的生活与行为方式的人，有能力将世上残余的邪灵驱逐出去？除了我们的救主，还有谁教导祂的门徒献上那些无血的、理性的祭献，这些祭献是通过祈祷和对天主隐秘的朝拜而进行的？\par

10. 因此，在普世之下，祭台被建立，教堂被奉献，在其中，这些属灵的、理性的祭献作为神圣的敬礼，由每个民族献给独一的至高天主。再者，除了祂，还有谁以无形而隐秘的力量，压制并彻底废除了那些以火和烟献上的血腥祭献，以及残忍无知的活人祭献？这一事实甚至由外教历史学家本身所证实。因为直到救主神圣教义公布之后，大约在哈德良统治时期，活人祭献的习俗才被普遍废止。\par

11. 如此明显的是我们救主死后权能与能力的证明。谁能找到如此顽固的灵魂，以致拒绝赞同真理，不肯承认祂的生命是神圣的？我所描述的这些行为是由活人所行，而非死人；可见的行为对我们来说，是那些不可见之事的证据。仿佛就在昨天，不虔敬和无神之徒扰乱并破坏了人类社会的和平，并拥有强大的力量。但这些，一旦生命离去，便匍匐在地，如同粪土一般毫无价值，无气息，无动作，失去言语，身后没有留下名声或纪念。因为这就是死人的状况；那不再活着的人就是虚无：一个虚无的人如何能做出任何行为呢？然而，祂的存在又怎能被质疑呢？祂的积极能力和能量甚至比那些仍然活着的人还要强大。虽然祂对肉眼不可见，但辨识能力并不在外在感官上。我们不通过身体感觉来理解艺术的规则或科学的理论；也没有任何眼睛能辨识人的心灵。更何况是天主的能力呢？在这种情况下，我们的判断是从显而易见的结果形成的。\par

12. 正因如此，我们必须判断我们救主的无形能力，并藉其显明的效果来决定：我们是应当承认祂现在仍在进行的大能作为是一位活生生的主体所行的；还是应将它们归给一个不存在者；或者，最后，是否这探询本身是荒谬和不一致的。因为我们有什么理由能断言一个不存在者的存在呢？既然所有人都承认，不存在的东西缺乏那种能力、能量和行动，因为这些是活物的特征，而相反则是死物的特征。\par

% structure: p0212 source=h2 target=section reason=relative-heading-rank; base=h2

\section{第十七章}

1. 如今时候已到，我们当思考我们的救主在我们这时代的作为，并默观永生天主的活泼行动。因为这些伟大的作为，除了作为一位真正享有天主生命的活泼行动者的生动明证，我们还能如何描述呢？若有人询问这些作为的性质，现在就请他留心。\par

2. 不久前，有一类人，受狂热驱使，又有同等权势和军事力量支持，显明他们对天主的敌意，摧毁他的教会，将敬拜他的建筑夷为平地。总之，他们用各种方式攻击那看不见的天主，用无数亵渎言语的箭矢袭击他。但那看不见者以看不见的手施行了报复。\par

3. 他仅凭意愿的一个命令，他的仇敌就全然毁灭。他们不久前还在极大的亨通中昌盛，被同胞高举，视为堪受神圣尊荣，并享有持续的权势和荣耀，只要他们与后来所反对的那位保持和平与友善。然而，一旦他们胆敢公开反抗他的旨意，列阵以他们的神对抗我们所敬拜的那位，立刻，按照他们所攻击的那天主的旨意和权能，他们所有人都受到与其胆大妄为相符的审判。他们被迫屈服并逃避他的权能，一同承认他的神性，并急忙撤销先前尝试的措施。\par

4. 因此，我们的救主毫不迟延地到处树立这胜利的纪念物，再次以圣殿和祝圣的祈祷之所装饰世界。在每座城市和村庄，不，在所有国家，甚至在蛮荒之地，他命令建立教会和神圣建筑，以尊崇至高的天主和万有的主。因此，这些祝圣的建筑物被认为配得上承继他的名，不是从人得名，而是从主自己得名，由此它们被称为教会（或主的房屋）。\par

5. 现在，让任何人出来告诉我们，在如此彻底的荒凉之后，是谁从地基到屋顶恢复了这些神圣建筑？当一切希望似乎消亡时，是谁使它们以比从前更宏伟的规模崛起？这确实令人惊叹：这重建并非在那些天主的仇敌死后，而是在这些建筑的摧毁者仍在世时发生，以至于他们以他们自己的言语和法令撤销了他们的恶行。他们这样做，并非处在兴盛与安逸的阳光之下（因为那时我们或许会以为仁慈或宽容是原因），而是在遭受天主惩罚打击之时。\par

6. 再者，在如此多次逼迫的风暴之后，甚至就在危险关头，是谁能够使普世如此多献身于哲学和事奉天主的人，以及那些献身于终身贞洁与纯净生活的圣洁童女团体，仍服从于他天上的诫命？是谁教导他们喜乐地坚持长久的禁食，并拥抱严格而一致的克己生活？是谁说服众多男女投身于神圣事物的研习，宁愿选择适于理性灵魂的灵性食物，而非身体的滋养？是谁教导野蛮人和农夫，甚至柔弱的妇女、奴隶和儿童，简言之，所有民族的无数群众，以轻看死亡而生活？他们相信灵魂的不朽，意识到人类的行为被那无误的正义之眼所注视，期待天主对义人和恶人的报应，因此忠于正义和美德的生活。否则，他们不能坚持虔诚的道路。确实，这些是现在我们的救主独自——也唯独他——所行的作为。\par

7. 现在，让我们从这些主题转而通过下面这样的探究，试图说服反对者顽固的心智。那么，无论你是谁，请上前来，说出理性的话语：开口，不是出于无知的心，而是出于智慧和光照的头脑：我说，在与你自己的深切严肃交谈之后，请开口。在那些至今闻名于世的圣贤中，有谁像我们的救主那样，从远古时代就被先知预言所预知和宣告给那曾经蒙天主眷顾的希伯来民族？但他的出生地、他来临的时期、他的生活方式、他的奇迹、言语和伟大作为，都已被这些先知的圣卷所预言和记载。\par

8. 再者，有谁如此现时地报复针对自己的罪行，以至于作为他们不敬的直接后果，整个犹太民族被无形的力量分散，他们的王座被彻底移除，他们的圣殿及其圣物被夷为平地？除了我们的救主，有谁既预言了那邪恶民族的事，又预言了他教会普世的建立，并同样以事实证实了二者？关于这些恶人的圣殿，我们的救主说：\Quote{你们的家要成为荒场留给你们} \BibleRef{玛 23:38}，又说：\Quote{在这里决不留一块石头在另一块石头上而不被拆毁的}。关于他的教会，他说：\Quote{我要在这磐石上建立我的教会，阴间的门决不能战胜它} \BibleRef{玛 16:18}。\par

9. 那召选无名无学的渔夫，使他们成为人类立法者和导师的权能，是何等奇妙！从他应许要使他们成为渔人的渔夫，并如此成就的应许中，我们找到他神性何等清楚的证明！他所赐予的权能和德能，使他们写下并出版了如此权威的著作，以至于被翻译成每一种文明和野蛮的语言，被所有民族诵读思考，其中的教义被当作天主的圣言而受到信任！\par

10. 他对未来的预言是何等奇妙，他的门徒们蒙受警告，说他们将被带到君王和统治者面前，将遭受最严厉的惩罚，并非因犯罪，而仅仅因宣认他的名，这见证何其伟大！或者，谁能充分描述那力量，他以这力量使他们预备好怀着甘愿的心去受苦，并使他们装备以虔敬的盔甲，在斗争中保持不屈不挠的精神恒心？\par

11. 或者，我们如何能足够钦佩那灵魂的坚定刚毅，它不仅坚固了他的直接追随者，也坚固了他们的后继者，直到我们现今的时代，使他们为证明对至高天主的忠忱而喜乐地忍受每一种刑罚和每一种折磨？再者，有哪位君王能将自己的统治延续如此漫长的时代？还有谁有能力在死后作战，战胜每一个敌人，征服每一个野蛮和文明的民族和城市，并以不可见和隐秘的手制服他的对手？\par

12. 最后，也是最重要的，哪一张诽谤的口敢质疑我们已提及的、由他的能力在世界各地建立的普世和平？因为所有国家之间的相互和睦与和谐，在时间上与我们的救主之道和宣讲在全世界的扩展相吻合：这是天主诸先知在远古时代所预言的诸事件的会合。恩慈的皇帝啊，若我试图只用一个视角展示我们救主之神圣权能的那些确凿证据（这些证据甚至在今天也可见其效果），白昼本身也将不够我用；因为在文明或野蛮民族中，从来没有一个人曾显示出像我们的救主那样的神圣德能。\par

13. 但我为何谈论人呢？因为所有民族视为神祇的存在中，没有一个曾显现在地上拥有像他那样的权能。如果有，那就让事实现在被证明。哲学家们，请你们前来告诉我们，有哪一位神祇或英雄曾为人所知，像我们的救主那样传授了永生和天国的道理？有谁像他那样说服了世界各地的群众追求神圣智慧的准则，将希望寄托于天堂本身，并期待为爱天主之人所预备的住所？有哪一位以人形显现的神祇或英雄曾从日出之地到日落之地（几乎与太阳之光同等的旅程）行进，并用他的教义那光明荣耀的光辉照亮人类，使地球上的每一个民族都联合敬拜独一的真天主？又有哪一位神祇或英雄曾像他那样，废除了文明或野蛮民族中的一切神祇和英雄；命令所有人不得接受神圣的尊荣，并要求服从那命令；然后，尽管独自与所有权力冲突，却彻底摧毁了敌对的军队；战胜了历代的神祇和英雄，使自己在普天下的每一区域都被万民承认为天主的独生子？\par

14. 还有谁曾命令居住在这广袤地球上的大陆和岛屿上的万民每周在主日聚集，并遵守它作为节日——不是为了满足身体，而是为了通过神圣真理的教导来坚强灵魂？有哪一位神祇或英雄，像我们的救主那样遭受如此激烈的冲突，却竖立了战胜每一个敌人的胜利纪念碑？因为他们确实从始至终不断地攻击他的教义和他的子民；但他，那不可见的一位，凭借隐秘力量的运用，提升了他的仆人和他们敬拜的神圣殿宇至荣耀的高峰。\par

但我们为何还要徒然试图详细描述那些我们救主权能的神圣证据？这些证据没有任何语言可以充分表达；它们本身无需我们的话语，却以最响亮的声音向那些心灵之耳向真理敞开的人呼吁。确实，这是一个奇怪的、奇妙的事实，在人类生活的编年史中无与伦比：我们所描述的福惠竟被赐予必死的我们，而那位实在是惟一的、永恒的天主圣子，竟如此可见于地上。\par

% structure: p0228 source=h2 target=section reason=relative-heading-rank; base=h2

\section{第十八章}

1. 然而，仁慈的君主，这些话在你听来或许显得多余，因为你已藉着屡次亲身的经历，确知我们救主的天主性；你更以行动而非言语，向全人类宣讲了真理。你或许会在闲暇之时，屈尊为我们讲述你的救主所赐予你的丰盛显现，以及他在你睡眠时多次显现给你的异象。我并非指那些向我们隐藏的隐秘提示，而是指他灌输到你心灵中、且关乎全人类普遍利益与福祉的法则。你将亲自以恰当的言辞，述说你的神圣盾牌与护卫者在战斗中赐予你的可见保护，你公开与隐秘仇敌的覆灭，以及他在危难时刻的及时援助。你要将困厄中的解脱、独处时的护卫、绝境中的对策、对未来的预知、你对公共利益的远见、你探究疑难问题的能力、你对重大事业的领导、你对政务的管理、你的军事部署和对各部门弊端的纠正、你对公共权利的条例，以及最后，你为众人共同福祉所制定的法律，都归于他。你或许还会向我们详述那些对我们隐秘、惟独你知晓、并珍藏于你王室记忆如同秘密宝库中的恩宠细节。毫无疑问，正是这些理由，这些令你心悦诚服的救主大能的证据，使你建造了那座神圣的建筑，向所有信徒与非信徒展示他战胜死亡的凯旋纪念碑，一座献给圣洁天主的圣洁殿宇；使你奉献了那些崇高而壮丽的不朽生命与天国的纪念物；使你呈献了纪念全能救主胜利的礼品，这些礼品充分彰显了奉献者——你——的帝王尊严。你用这样的纪念物装饰了那座见证永生的建筑；如此，仿佛以帝王的文字，将胜利与凯旋归于天主的圣言；如此，以清晰无误的声音，以言行向万国宣告你对他圣名的虔诚宣认。\par

\SourceRecord{Translated by Ernest Cushing Richardson.；Nicene and Post-Nicene Fathers, Second Series；Vol. 1.；Edited by Philip Schaff and Henry Wace.；Buffalo, NY: Christian Literature Publishing Co.,；1890.；<http://www.newadvent.org/fathers/2504.htm>.}
